\subsection{One exclusion criterion for several geometric cases}
\label{sec:gradientexclusion}

The smooth case does not need a list of critical-point formulas.
The same argument also treats many singular leading forms. Its
hypotheses must hold in every complex direction. A bound only on
the original real integration space would not suffice.

We state the criterion for a polynomial exponential amplitude, then
specialize it to the Gaussian. Norms are Hermitian norms on complex
affine space; the Gaussian quadratic remains holomorphic.

\begin{theorem}[A gradient criterion]\label{thm:gradientcriterion}
Let $P,q\in\C[x_1,\ldots,x_n]$, with $P$ nonconstant and
$\deg q\leq J$, where $J\geq1$. Suppose that outside a compact set
\begin{equation}\label{eq:gradientconditions}
 \|\nabla P(x)\|\geq c\|x\|^{J-1},\qquad
 \frac{\|\operatorname{Hess}P(x)\|}{\|\nabla P(x)\|^2}\leq C
\end{equation}
for some $c>0$ and $C<\infty$. For every fixed polynomial $A$,
no nonzero flat period of
\[
 A(x)e^{tP(x)-q(x)}\,dx_1\cdots dx_n
\]
is a finite exponential-polynomial sum
$\sum_\lambda e^{\lambda t}p_\lambda(t)$ with polynomial
$p_\lambda$. In particular it cannot be a nonzero constant.
\end{theorem}
\begin{proof}
We first justify the full cycle space and the tails. The gradient is
nonzero outside a compact set, so the affine critical set is compact
and algebraic, hence finite. The gradient bound is stronger than the
condition excluding vanishing cycles at infinity. The finite-critical
thimble decomposition therefore applies, with the whole local Milnor
homology at each isolated critical point, not just its Morse
approximation~\cite[Part~I, Sections~1--4]{Pham}.
When several critical points have the same value, take disjoint
Milnor balls and apply excision to their complement. This gives the
direct sum of their full local homology before any values are grouped.

Here the Gaussian or polynomial amplitude can be kept on these fixed
$P$-cycles. Choose a generic argument $\theta$ of $t$. Outside the
local critical neighbourhoods, lift a descending phase ray by
\[
 P(x)=\lambda-e^{-i\theta}s,\qquad
 \frac{dx}{ds}=-e^{-i\theta}
                    \frac{\overline{\nabla P(x)}}{\|\nabla P(x)\|^2}.
\]
For large $s$ the lift is outside the compact set, because its phase
value is unbounded. The first condition in
\eqref{eq:gradientconditions} gives
$\|x\|^J\leq C_1(1+s)$. The real differential of this vector field
has norm at most a constant times the second ratio in
\eqref{eq:gradientconditions}, hence is bounded. Transported
$(n-1)$-cycle volumes grow at most as $\exp(C_2s)$; the radial speed
is bounded as well. Also
$|A(x)e^{-q(x)}|\leq C_3(1+s)^M\exp(C_4s)$ for some $M$.
Thus $\exp(-|t|s)$ dominates the complete tail, uniformly for large
$|t|$ in a smaller sector. The same estimate applies to every fixed
polynomial amplitude.

We must also check that these are a full basis for the actual phase,
rather than selected admissible cycles. Put
\[
 F_\sigma=P-\sigma q/t,\qquad 0\leq\sigma\leq1.
\]
Since $\|\nabla q\|\leq C_5\|\nabla P\|$ outside the same compact
set, all large $|t|$ give
\[
 \|\nabla F_\sigma\|\geq\tfrac12\|\nabla P\|.
\]
The ratio involving the Hessian stays uniformly bounded too. Indeed
$\|\operatorname{Hess}q\|=O(\|x\|^{J-2})$ for $J\geq2$, while
that Hessian is zero for $J=1$.
All critical points and their values therefore remain in fixed compact
sets. On the preimage of a phase half-plane beyond those values, the
parameter field
\[
 W_\sigma=\frac{q}{t}
       \frac{\overline{\nabla F_\sigma}}{\|\nabla F_\sigma\|^2}
\]
satisfies $\partial_\sigma F_\sigma+dF_\sigma(W_\sigma)=0$ and
$\|W_\sigma\|\leq C_6(1+\|x\|)/|t|$. Its flow is complete for
$0\leq\sigma\leq1$. Extend it by a cutoff away from that half-plane.
It gives an isotopy of the ambient relative pairs defining the full
rapid-decay homology. This explicitly prevents an infinity class
from appearing during the deformation.

The fixed $P$-thimbles have uniformly descending tails for every
$F_\sigma$, by the preceding estimates. Truncate them beyond a common
large phase value. Their boundaries stay in the descending half-plane
throughout the isotopy, so they define its flat relative classes.
They are a full basis at $\sigma=0$ and remain one at $\sigma=1$.
Finite critical points may merge, and their critical values may
coincide; the argument uses the full relative pairs and does not
require a Morse family. The same cycles are flat as $t$ varies in
the chosen sector.

These improper chains and their homotopies are admissible for the
period pairing. One can integrate first on truncated chains and then
pass to the limit: the tail bound makes every boundary term tend to
zero. Indeed the differential of $W_\sigma$ grows at most as
$C(1+\|x\|^J)/|t|$; along the controlled tails this is
$C'(1+s)/|t|$. The homotopy volumes have at most this exponential
growth, still dominated by the descending factor for all large $|t|$.
Thus Stokes and differentiation under the integral hold throughout;
the relative-pair calculation does not depend on a convergence claim
for the original real Gaussian integral.

It remains to determine the local asymptotics at a finite critical
point with value $\lambda$. For an analytic amplitude germ, the
vanishing-cycle density is a finite sum of terms
$s^{\alpha+k}(\log s)^j$, with analytic coefficient series and
$\alpha>-1$. Its Laplace asymptotics consequently have the form
\[
 e^{\lambda t}\sum_{\alpha,k,j}
       c_{\alpha,k,j}\,t^{-\alpha-k-1}(\log t)^j,
 \qquad k\geq0.
\]
This is the local analytic result of
Pham~\cite[Part~II, Section~1.1, (1.5)--(1.6) and the footnote on
printed page~26]{Pham}; it permits degenerate isolated critical
points. Pham credits Maisonobe and Rombaldi for the local solution
comparison used there. The amplitude $A e^{-q}$ is analytic. Every nonzero
term has a strictly negative power of $t$ and tends to zero after
removing $e^{\lambda t}$. The tails outside the local neighbourhood
are exponentially negligible by the uniform bound already proved.

The finite-critical decomposition accounts for the full module.
Group equal critical values before separating the distinct formal
exponential factors. Within each such factor the displayed local
expansions contain only decaying powers and logarithms. A nonzero
polynomial prefactor cannot occur. The zero function is allowed when
all these contributions vanish. This proves the theorem.
\end{proof}

\begin{corollary}[Original Gaussian markers]\label{cor:gradientmarkers}
For the standard Gaussian quadratic $q=\sum x_i^2/2$, suppose
\eqref{eq:gradientconditions} holds with $J=2$. Then every original
marked series $\Phi_A$ is either zero or is not a finite
exponential-polynomial sum. In particular $P$ is neither pure nor safe.
\end{corollary}
\begin{proof}
The exact top de Rham quotient and connection are
\[
 M=\C(t)[x]\Big/\sum_i(\partial_i+tP_i-x_i)\C(t)[x],
 \qquad\nabla[A]=[A_t+PA].
\]
Gaussian integration by parts makes $[A]\mapsto\Phi_A$ horizontal
on this vector-space quotient. The spatial and parameter derivatives
commute; no ordinary Jacobian-ideal quotient is substituted.
If some $\Phi_A$ is a
nonzero finite exponential-polynomial sum, a constant-coefficient
differential operator isolates one exponential and kills its
polynomial prefactor, giving $c e^{\beta t}$, $c\ne0$.
Because
\[
 \partial_t\Phi_A=\Phi_{PA},
\]
this is still an original fixed polynomial marker. Replace $P$ by
$P-\beta$ and divide the marker by $c$. Its formal seed value is one.
Lemma~\ref{lem:seedfunctional} gives an actual constant period for
that marker, contradicting Theorem~\ref{thm:gradientcriterion}.
Its comparison hypotheses hold with base $t$ and phase $tP-q$ on
$\mathbb A^1\times\mathbb A^n$, after removing the finite exceptional
base set. No convergence of the original Gaussian series is assumed.
Finally $\Phi_1(0)=1$, so it is a nonzero series; purity or safety
would make it a polynomial.
\end{proof}

\begin{corollary}[Smooth leading form, in every degree]
\label{cor:smoothleading}
If $\deg P=D\geq2$ and the homogeneous leading form satisfies
$\nabla F_D(v)=0$ only for $v=0$, then $P$ is not pure. In fact
Corollary~\ref{cor:gradientmarkers} holds for every original marker.
\end{corollary}
\begin{proof}
Compactness of the complex unit sphere gives
$\|\nabla P(x)\|\geq c\|x\|^{D-1}$ for large $\|x\|$.
Also $\|\operatorname{Hess}P\|=O(\|x\|^{D-2})$, so its ratio to
the squared gradient is $O(\|x\|^{-D})$. Both conditions for $J=2$
follow, with all lower terms retained.
\end{proof}

\begin{lemma}[The ordinary-node test]\label{lem:nodecoercivity}
Let $P=F_3+F_2+F_1+F_0$ in $n\geq2$ variables. Suppose its
projective singular set is finite and every singular point $[v]$
has $\ker\operatorname{Hess}F_3(v)=\C v$. If $F_2(v)\ne0$ at
every such point, then
$\|\nabla P(x)\|\geq c\|x\|$ outside a compact set. Consequently
$P$ is not pure under the Gaussian functional.
\end{lemma}
\begin{proof}
Otherwise choose an unbounded sequence with
$\|\nabla P(x)\|/\|x\|\longrightarrow0$. A projective limit $[v]$
is singular for $F_3$. Choose a fixed complement to $\C v$ and write
$x=sv+y$, where $y=o(s)$. With $H=\operatorname{Hess}F_3(v)$,
homogeneity gives the exact expansion
\[
 \nabla P(sv+y)=sHy+\nabla F_3(y)+s\nabla F_2(v)
                              +\nabla F_2(y)+\nabla F_1.
\]
The map $H$ is injective on the chosen complement. Divide by $s$.
The error from $\nabla F_3(y)/s$ is $o(\|y\|)$, so the equation
first forces $y$ to stay bounded. A limit $y_0$ satisfies
$Hy_0+\nabla F_2(v)=0$. Pair bilinearly with $v$. Symmetry and
$Hv=0$ kill the first term, whereas Euler's identity gives
$v\cdot\nabla F_2(v)=2F_2(v)$, a contradiction.
The Hessian of a cubic is $O(1+\|x\|)$, so the second condition
in \eqref{eq:gradientconditions} follows from the first. Apply
Corollary~\ref{cor:gradientmarkers}.
\end{proof}

\subsection{A common residual family for ordinary nodes}
\label{sec:ordinaryresidual}

For a pure cubic with only ordinary leading nodes, the preceding
lemma forces $F_2$ to vanish at some node. An isotropic such node
gives the frame already proved in Lemma~\ref{lem:frame}. Every
nonisotropic such node has the following single normal form.

\begin{lemma}\label{lem:ordinarynormal}
Let $[w]$ be an ordinary node of a ternary $F_3$, with $F_2(w)=0$
and $w^Tw\ne0$. An invertible affine change of coordinates puts
the pair consisting of $P$ and the Gaussian quadratic into
\begin{align*}
 P={}&T(uv+h)+a u^3+d v^3+b u^2+c v^2+e u+f v+k,\\
 q={}&\tfrac12(Au^2+Bv^2+mT^2+2Cuv+2DuT+2EvT)
                  +\ell_u u+\ell_v v+\ell_T T+q_0,
\end{align*}
where $m\ne0$ and $\det\operatorname{Hess}q\ne0$.
All Gaussian means and cross terms must be retained. Its leading
curve is irreducible nodal when $ad\ne0$, a transverse line and
smooth conic when exactly one of $a,d$ vanishes, and three
nonconcurrent lines when $a=d=0$.
\end{lemma}
\begin{proof}
Take $T$ in the direction $w$ and choose a complementary plane.
The singular-point equations make $F_3$ linear in $T$; its
coefficient is a nondegenerate binary quadratic because the node
is ordinary. Factor that quadratic and choose its two linear
factors as coordinates $u,v$. Thus $F_3=Tuv+B_3(u,v)$.
Replace $T$ by $T$ plus a linear form in $u,v$ to remove the
$u^2v$ and $uv^2$ terms of $B_3$.

There is no $T^2$ term in $F_2$, precisely because $F_2(w)=0$.
The full coefficient of $T$ is therefore
$uv+\alpha u+\beta v+h_0$. Translations of $u,v$ make it $uv+h$.
The cubic terms outside $T$ remain pure cubes. A constant shift
of $T$ removes the remaining quadratic $uv$ term. This gives the
displayed $P$. The direction of $T$ has not changed, so $m\ne0$;
invertibility of the affine change preserves the nondegenerate
Gaussian Hessian, while translations and the shifts of $T$ produce
the means and cross terms shown.

If $ad\ne0$, the polynomial $Tuv+a u^3+d v^3$ is irreducible:
an independent factor of a polynomial linear in $T$ would divide
both $uv$ and $a u^3+d v^3$, whose gcd is one. Its gradient shows
that its sole singular point is the stated ordinary node. If
$a=0,d\ne0$, it is $v(Tu+d v^2)$, with a smooth conic meeting
the line transversely at two points; the other one-zero case is
symmetric. If both vanish, it is the coordinate triangle $Tuv$.
\end{proof}

This consolidates three leading-curve rows into one residual family.
It is a reduction, not yet an exclusion of that family. In particular
the full quadratic $q$, including its means, cannot be replaced by
a convenient diagonal form during the remaining period argument.

\subsection{Two controls on the scope of the criterion}

Ordinary tameness of $P$ alone is insufficient for the Gaussian
conclusion. Let $P=\ell^Tx$, with $\ell=(1,i)/\sqrt2$ in two
variables. Its gradient is constant and nonzero, but
$\Phi_1=\exp(t^2\ell^T\ell/2)=1$. The phase $tP-q$ has its
critical point at $x=t\ell$, escaping with action zero. Thus an
isolated critical set for $P$, or the weaker gradient bound away
from zero, would lose exactly the contribution that purity needs.
Condition~\eqref{eq:gradientconditions} with $J=2$ prevents this.

The distinction between an original marker and a time-dependent
amplitude also remains essential. For
$P=(X_1^2+X_2^2)/2$ one has $\Phi_1=(1-t)^{-1}$ and
$(1-t)\Phi_1=1$. This phase satisfies the gradient criterion.
The unit in its rational differential module is represented by a
time-dependent amplitude. No original fixed polynomial marker has
a nonzero polynomial period, by Corollary~\ref{cor:gradientmarkers}.
The abstract module alone does not record that distinction.

\subsection{Ordinary nodes with a finite affine critical set}
\label{sec:ordinaryfinite}

The three ordinary-node families admit one argument whenever the
critical set of $P$ is finite. In particular we need not solve for
each escaping critical point, assume it is Morse, or separate its
action from the actions of other escaping points. The Gaussian
quadratic itself detects the unwanted contribution at infinity.

Retain the notation and the full pair $(P,q)$ of
Lemma~\ref{lem:ordinarynormal}. Write $q_2$ for the homogeneous
quadratic part of $q$. The absence of a Gaussian frame is exactly
the following condition at the leading nodes:
\begin{equation}\label{eq:nodenoframe}
 F_2(w)=0\quad\Longrightarrow\quad q_2(w)\ne0.
\end{equation}
Indeed the Gaussian bilinear form in these affine coordinates is
$2q_2$; the value of $F_2$ at a leading singular direction is
unchanged by translation. In the common normal form the $T$ node
already satisfies this condition because $m\ne0$. The only extra
conditions are $c\ne0$ or $B\ne0$ when $d=0$, and $b\ne0$ or
$A\ne0$ when $a=0$.

\begin{lemma}[The stationary action identity]\label{lem:actionidentity}
Put $r=1/t$, and let $x(r)$ be an algebraic critical branch of
$S=P/r-q$. A prime means total differentiation along the branch.
Then
\begin{equation}\label{eq:actionidentity}
 S'=-P(x(r))/r^2,\qquad q(x(r))=-S-rS'.
\end{equation}
If $x(r)$ is bounded, the polar part of $S$ is $\alpha/r$, where
$\alpha=P(x(0))$. If it escapes and \eqref{eq:nodenoframe} holds,
its polar part is neither zero nor $\alpha/r$ for any constant
$\alpha$.
\end{lemma}
\begin{proof}
Use a finite Puiseux cover so every branch is meromorphic. At a
critical point the chain-rule term $\nabla_x S\cdot x'$ vanishes,
giving both identities. For a monomial $c r^{-\kappa}$ of the
polar part of $S$, its contribution to $-S-rS'$ is
$(\kappa-1)c r^{-\kappa}$. Thus bounded $q(x(r))$ permits only
the pole $\alpha/r$ in $S$. Its coefficient is $P(x(0))$ when
$x$ is bounded.

For an escaping branch let $R=\|x\|$. The critical equation is
$\nabla P=r\nabla q$, so $\|\nabla P\|/R\longrightarrow0$.
The leading projective direction is a singular point of $F_3$.
The ordinary-node argument in the proof of the preceding coercive
lemma works at this one direction: write $x=sw+y$, $y=o(s)$,
in a fixed complementary plane. The injectivity there of
$H=\operatorname{Hess}F_3(w)$ forces $y$ bounded from
\[
 \nabla P/s=Hy+\nabla F_2(w)+O((1+\|y\|)^2/|s|)=o(1).
\]
Pairing with $w$ then gives $F_2(w)=0$. Condition
\eqref{eq:nodenoframe} gives $q_2(w)\ne0$, so $q(x(r))$ has a
pole: its leading quadratic term cannot be cancelled by its
lower-degree terms. Identity~\eqref{eq:actionidentity} now rules
out both claimed polar parts. This argument includes all escape
scales and all coincident actions.
\end{proof}

We specify the classical rank input before using it. This is a
convergent algebraic application of formal stationary phase, not
an extension to arbitrary nonconvergent coefficients or rings.

\begin{lemma}[Rank accounting for a tame algebraic pencil]
\label{lem:tamepencilrank}
For the present no-frame pencil, shrink a punctured $r$ disk to
avoid its finitely many exceptional parameters. Its top twisted
de Rham module has rank
\[
 N=\sum_{\text{critical points of }S_r}\mu.
\]
After a finite cover, the rank of its block with nonzero
exponential polar part $\gamma$ is the sum of these generic local
Milnor numbers over critical branches with action polar part
$\gamma$. Equal actions add their multiplicities. Consequently
the rank of the zero-exponential block is the sum of the generic
Milnor numbers of the branches with bounded action.
\end{lemma}
\begin{proof}
For fixed small nonzero $r$, the homogeneous quadratic term of
$S_r$ at a leading node $w$ is $F_2(w)/r-q_2(w)$. It is nonzero
at every node, by \eqref{eq:nodenoframe} and exclusion of finitely
many parameters. The ordinary-node coercivity lemma gives
$\|\nabla S_r\|\ge c_r\|x\|$ outside a compact set. The
critical locus is therefore finite: a positive-dimensional affine
algebraic component would be unbounded. Pham's tame Milnor
decomposition~\cite[Part~I, Sections~1--4]{Pham} gives the displayed
rank, retaining the Milnor number of every degenerate point.

Here are the checks needed to apply the direct-image rank formula.
Take $U=\C_r^*\times\C_x^3$, $f=r$, $g=S_r$, and the regular
algebraic module $\mathcal O_U$. Form the ordinary direct image
by $(f,g)$, followed by exponential integration in the $g$
coordinate. All these are algebraic maps; $g$ is regular on $U$.
The ordinary direct image is regular holonomic. Over the punctured
$r$ disk its relevant cohomology is the middle, reduced
Gauss--Manin module of the fibers of $S_r$. The other, constant
fiber cohomology has zero exponential integral in the affine $g$
line. Punctual parameter terms disappear on the punctured disk.

There are no extra finite-$g$ singular curves from spatial infinity.
The coercive estimate is locally uniform on a compact subset of
the punctured $r$ disk. For $g$ in a compact set, the identity
$P=r(g+q)$ on a fiber gives
\[
 \partial_r S_r=-(g+q)/r=O(1+\|x\|^2).
\]
Lifting $r$ variation by the normalized gradient therefore uses a
vector field of at most linear growth. Lifting $g$ variation uses
a field $O(1/\|x\|)$. These fields are complete outside a fixed
ball on each such parameter neighborhood. They trivialize the
outer fiber pairs, so all nonvertical finite-$g$ discriminant
curves come from affine critical values. Along such a curve the
conormal multiplicity of the middle Gauss--Manin module is the
sum of the local Milnor numbers: a transverse target disk computes
precisely the corresponding isolated vanishing cycles. This also
applies to persistent degenerate points and equal critical values.

First make a common finite Puiseux base change of all the
critical-value curves. Apply the nonramified case of Roucairol's
Theorem~0.1(2), proved in her Section~1.2, and Theorem~0.3
\cite{RoucairolFormal} to this identified ordinary direct image.
They give exactly the stated nonzero polar parts and ranks.
Count every graph after this base change, including distinct
conjugates with the same polar part. For example the two graphs
$g=r^{-1}\pm r^{1/2}$ must contribute twice after $r=\rho^2$;
retaining just their one polar expression would lose rank.
Their rank statement requires no separation of the constant terms
of equal polar parts. The additional separation hypothesis in her
Theorem~0.2 concerns its monodromy formula, which we do not use.
Vertical components at $r=0$ and the boundary $g=\infty$ do not
add polar branches in Theorem~0.1. Subtracting the nonzero-block
ranks from $N$ proves the last assertion. In particular we have
not inferred the regular rank from an Euler characteristic alone.
\end{proof}

\begin{proposition}\label{prop:ordinaryfinite}
Suppose the common ordinary-node pair satisfies
\eqref{eq:nodenoframe} and $P$ has finite affine critical set.
Then $P$ is not pure. More generally every original Gaussian
marked period is either zero or is not a finite sum of
polynomials in $t$ times exponentials $e^{\beta t}$.
\end{proposition}
\begin{proof}
Every bounded critical branch of $S_r$ limits to a critical point
of $P$. Disjoint small Milnor neighborhoods of these finitely many
points contain, for the deformation $P-rq$, total critical
multiplicity $\mu_p$ in the neighborhood of $p$. This is
conservation of the local intersection multiplicity of the
gradient, and does not require a Morse perturbation. By
Lemma~\ref{lem:actionidentity} its action has polar part
$P(p)/r$. Every escaping branch has nonzero polar part. Hence
Lemma~\ref{lem:tamepencilrank} identifies the entire zero-exponential
rank as
\begin{equation}\label{eq:ordinaryzerorank}
 \sum_{p:\,dP(p)=0,\,P(p)=0}\mu_p.
\end{equation}

We construct that many independent local formal functionals and
compute their monodromy over the unramified field $\C((r))$.
At an isolated critical point at least one of $u,v$ is nonzero.
Indeed if both vanish, criticality requires $h=e=f=0$, in which
case the whole $T$ axis is critical. In the chart $u\ne0$ put
\[
 s=uv+h,\quad R(u)=G(u,-h/u),\quad
 T'=T+\frac{G(u,(s-h)/u)-G(u,-h/u)}s,
\]
where $G=P-T(uv+h)$. The quotient is analytic, and exactly
$P=sT'+R(u)$. The coordinate Jacobian is $u\ne0$. At a critical
point $p$ the one-variable critical multiplicity of $R$ is
$\mu_p$. An analytic one-variable coordinate $y$, with
$w=T'$, therefore gives the exact identity
\[
 P=\alpha+sw+y^k,\qquad k=\mu_p+1,\quad\alpha=P(p).
\]
The volume Jacobian and $e^{-q}$ multiply every amplitude by an
analytic unit $V(y,s,w)$; no Gaussian mean or cross term has been
discarded.

For $j=1,\ldots,k-1$ define a formal local functional on an
analytic monomial $y^a s^b w^c$ as zero unless $b=c$ and
$a=j-1+\ell k$ for an integer $\ell\ge0$. Otherwise its value is
\begin{equation}\label{eq:ordinarycharacters}
 e^{\alpha/r}r^{1+j/k}
       (-1)^{b+\ell}b!\,(j/k)_\ell\,r^{b+\ell}.
\end{equation}
Evaluate a polynomial form after multiplying its transformed
amplitude by $V$. At each power of $r$ only finitely many Taylor
terms contribute. The recurrences
\[
 I_{a+k}=-\frac{a+1}{k}r I_a,\qquad
 J_{b+1,b+1}=-(b+1)rJ_{b,b}
\]
prove spatial integration by parts and parameter horizontality.
Thus these are functionals on the full twisted quotient, not on
an assumed cyclic image. They exist even if the corresponding
asymptotic series diverge; their input coefficients here are the
Taylor series of the stated analytic coordinates.

They are independent. Prescribe, at each $p$, the jets
$1,y,\ldots,y^{k-2}$ on the line $s=w=0$ and zero jets at the
other points, using polynomial representatives of finite local
jets and the Chinese remainder theorem. After dividing column
$j$ by $e^{\alpha/r}r^{1+j/k}$, the leading evaluation matrix is
the jet map multiplied by the analytic unit $V$. It is triangular
with nonzero diagonal. Its higher terms are in $r\C[[r]]$.
For $\alpha=0$ these independent characters fill exactly the rank
in \eqref{eq:ordinaryzerorank}; there is no unaccounted regular
piece. Their formal monodromy eigenvalues, or their inverses if
one uses the dual module, are
\[
 \exp(2\pi i j/k),\qquad 1\le j<k.
\]
None is one. This conclusion is over $\C((r))$. We do not replace
a noninteger residue by its integer multiple after ramification.
The unique zero-exponential projector commutes with descent, and
the displayed unramified coefficient matrices identify its whole
block. Every other block has a nonzero exponential polar part.
The ambient module therefore has no trivial composition factor.

For completeness this excludes the formal Gaussian unit without
choosing an actual analytic period. Formal Gaussian integration
annihilates $\partial_i+tP_i-q_i$ coefficientwise. If its value
on a fixed marker $A$ were a nonzero constant, restriction to the
rational differential submodule generated by $A$ would give a
surjection to the trivial module: apply the functional to every
rational relation between $A,\nabla_tA,\ldots$. A relation has
zero constant coefficient because all positive derivatives of
that constant vanish. Exact formal extension at infinity would
preserve this composition factor, contradicting what was proved.

Finally a nonzero finite exponential-polynomial can be reduced
to $c e^{\beta t}$ by a constant-coefficient differential
operator. On an original marker this operator replaces $A$ by
the fixed polynomial marker $T(P)A$. Replace $P$ by $P-\beta$
and divide the marker by $c$. The hypotheses and finite critical
set are unchanged, so the same contradiction applies. Taking
$A=1$ proves the purity exclusion.
\end{proof}

\subsection{An exhaustive critical-set split}

The finite-critical hypothesis in Proposition~\ref{prop:ordinaryfinite}
has a short exhaustive coefficient test. This prevents a generic
argument from silently discarding the positive-dimensional cases.

If $h\ne0$, the critical locus lies on $uv=-h$, and the equations
above identify it with $R'(u)=0$ in $\C^*$. The Laurent polynomial
\[
 R(u)=a u^3+b u^2+e u+k-fh/u+ch^2/u^2-dh^3/u^3
\]
is nonconstant unless $a=b=c=d=e=f=0$. A nonzero Laurent
derivative has only finitely many roots in $\C^*$. In the constant
case the critical set is the smooth conic $uv=-h$, $T=0$.

If $h=0$, outside the center the critical equations are
\[
 \begin{array}{lll}
 v=0,\ u\ne0:&3au^2+2bu+e=0,&T=-f/u,\\
 u=0,\ v\ne0:&3dv^2+2cv+f=0,&T=-e/v.
 \end{array}
\]
The center $u=v=0$ is critical exactly when $e=f=0$, and then
the whole $T$ axis is critical. Consequently the following table
is exhaustive and disjoint:
\begin{center}\small
\begin{tabular}{@{}p{0.51\textwidth}p{0.39\textwidth}@{}}
\toprule
Condition & Conclusion\\
\midrule
$h\ne0$, $(a,b,c,d,e,f)\ne0$ & Finite critical set; excluded above\\
$h\ne0$, $a=b=c=d=e=f=0$ & Critical conic; Corollary~\ref{cor:ordinaryconics}\\
$h=0$, $e=f=0$ & Center axes; Proposition~\ref{prop:axisexclusion}\\
$h=0$, $(e,f)\ne0$, $a=b=e=0$ & Critical hyperbola; Corollary~\ref{cor:ordinaryconics}\\
$h=0$, $(e,f)\ne0$, $d=c=f=0$ & Symmetric hyperbola; same corollary\\
$h=0$, $(e,f)\ne0$, $(a,b,e)\ne0$, $(d,c,f)\ne0$
 & Finite critical set; excluded above\\
\bottomrule
\end{tabular}
\end{center}
The two hyperbola rows cannot overlap because $(e,f)\ne0$.
All isolated degeneracies in the finite rows have already been
included through their Milnor multiplicities. The conic and hyperbola rows are proved in the next subsection.
The center-axis row is proved in Proposition~\ref{prop:axisexclusion},
using a full formal frame or the uniform actual-cut criterion.

\subsection{A critical conic in every polynomial degree}
\label{sec:hyperbolicsuspension}

The conic and hyperbola rows of the preceding table are manifestations
of the same calculation. It also works when one of the cubic terms is
replaced by a polynomial of arbitrary degree. We prove that stronger
statement to make clear what the mechanism actually uses.

\begin{theorem}\label{thm:hyperbolicsuspension}
Retain a full nondegenerate affine Gaussian quadratic
\[
 q=\tfrac12(Ax^2+By^2+mw^2+2Cxy+2Dxw+2Eyw)
                     +\ell_x x+\ell_y y+\ell_w w+q_0.
\]
Let $AB\ne0$ and
\[
 P=wxy+F(w)+k,\qquad F(0)=0,\quad F'(0)=h\ne0.
\]
Allow any polynomial degree $d_F=\deg F\ge1$, requiring $m\ne0$
when $d_F=1$. Then $P$ is not pure. Every original marked Gaussian
period is either zero or is not a finite exponential-polynomial.
The generic top twisted rank is $d_F+3$ for $d_F\ge2$ and five for
$d_F=1$. No Morse assumption on the Laurent phase obtained by
restricting $q$ to the conic, or on the isolated critical points of
$P$, is needed.
\end{theorem}

The nonzero quadratic poles at both ends of the conic are essential
to the rank-four construction below. The condition $AB\ne0$ states
this explicitly. Nondegeneracy of the full Gaussian alone would not
imply it. This theorem concerns complex polynomials and valued-field
formal frames; it does not assert an all-ring extension.

\begin{proof}
We first justify the full rank. Fix a generic $t\ne0$ and set
$S=tP-q$. Its gradient has a linear lower bound at complex spatial
infinity. To see this, suppose along an unbounded sequence
$\|\nabla S\|=o(R)$, where $R=\|(x,y,w)\|$. If $w$ stays
bounded, pass to a limit $w_0$ and normalize $(x,y)$ by $R$.
The first two gradient equations give
\[
 -AX+(tw_0-C)Y=0,\qquad (tw_0-C)X-BY=0.
\]
At least one of $X,Y$ is nonzero, and $AB\ne0$ makes both
nonzero. But the last gradient equation, divided by $R^2$, gives
$tXY=0$, a contradiction.

If $w$ is unbounded, the first two equations have a matrix with
inverse $O(1/|w|)$. Hence
$x,y=O(1)+o(R)/|w|$. First $R\asymp |w|$, then $x,y=O(1)$.
The last equation is therefore
\[
 tF'(w)-mw+O(1)=o(|w|).
\]
This is impossible for $d_F\ge3$, for $d_F=2$ with
$2t[w^2]F-m\ne0$, and for $d_F=1$ with $m\ne0$.
The same proof is uniform on compact subsets of these coefficient
conditions. Thus all phase-critical points stay in one ball during
any compact coefficient deformation inside that domain. They are
isolated, and their total Milnor number is constant by conservation
of local gradient-intersection multiplicity. Pham's tame Milnor
decomposition~\cite[Part~I, Sections~1--4]{Pham} identifies that
number with the full twisted rank.

For fixed $t$, the coefficient domain just used is the complement
of finitely many complex algebraic hypersurfaces, and is path
connected. We may compute the number at a generic diagonal
Gaussian with zero means, retaining $ABm\ne0$, and a generic $F$.
The critical equations there are
\[
 twy=Ax,\qquad twx=By,\qquad t(xy+F'(w))=mw.
\]
There are $d_F-1$ simple solutions with $x=y=0$ if $d_F\ge2$,
and one if $d_F=1$. The other solutions have
$w=\pm\sqrt{AB}/t$. At each of these two values, $x,y$ are
proportional and their nonzero product is prescribed by the last
equation; generically this gives two simple solutions. The counts
are $d_F+3$ and five, respectively. Uniform coercivity and the
connected deformation justify this specialization; a convenient
critical-point count by itself would not justify the rank for a
degenerate member.

We now construct a full formal frame. Set $r=t^{-1/2}$ and write
$F(w)=hw+bw^2+\sum_{j\ge3}f_jw^j$. Near the entire critical
conic $w=0$, $xy=-h$, use the exact coordinates
\[
 y=(rV-h)/x,\qquad w=rU.
\]
They give
\begin{align}
 tP-q&=k/r^2+UV+bU^2+\psi(x)+E_r(x,V,U),
                  &E_r\in r\C[x,x^{-1},V,U][r],\label{eq:conicphase}\\
 dx\,dy\,dw&=r^2\,\frac{dx}{x}\,dV\,dU,\label{eq:conicvolume}\\
 \psi(x)&=-\frac A2x^2-\frac{Bh^2}{2x^2}
                -\ell_x x+\frac{h\ell_y}{x}+Ch-q_0.\notag
\end{align}
All Gaussian cross terms and means are in $\psi$ or $E_r$.
The normal quadratic $UV+bU^2$ has Hessian determinant $-1$ for
every $b$. Its normalized formal Gaussian functional is well
defined and takes one to one.

The logarithmic Laurent quotient for $\psi$ has rank four, with
basis $1,x,x^2,x^3$. Here is a direct proof, including independence.
Its relation on $x^j$ is
\[
 jx^j-Ax^{j+2}-\ell_xx^{j+1}
             -h\ell_yx^{j-1}+Bh^2x^{j-2}=0.
\]
The top coefficient eliminates all exponents at least four without
creating a negative exponent; the bottom eliminates all negative
exponents without creating an exponent above three. If the image
of a nonzero Laurent polynomial under $x\partial_x+x\psi_x$
were supported in $[0,3]$, its highest exponent would be at most
one and its lowest at least two, which is impossible.

Let $\lambda_0,\ldots,\lambda_3$ be the dual basis on this
quotient, tensored with the normal Gaussian functional. For a
polynomial amplitude $a(x,y,w)$, define
\begin{equation}\label{eq:coniccharacters}
 L_i(a)=e^{k/r^2}r^2\lambda_i\!\left(
    e^{E_r}a(x,(rV-h)/x,rU)\right).
\end{equation}
Every coefficient is a finite Laurent-polynomial calculation, so
this is a well-defined formal series. Multiplication by $e^{E_r}$
conjugates the fixed twisted differential to the perturbed one.
Consequently these four functionals annihilate all spatial exact
forms. The moving volume is retained in the factor $r^2$.
For an explicit parameter check the coordinate velocity in the
original variables is
\[
 Z=\left(0,\frac{xy+h}{rx},\frac wr\right),
                   \qquad \operatorname{div}Z=2/r.
\]
Cartan's identity $\mathcal L_Z=d\iota_Z+\iota_Zd$ shows that
the pullback derivative differs from the original parameter
connection only by a spatial exact form and this volume
derivative. Thus \eqref{eq:coniccharacters} are horizontal for
the original full module. In particular, for every fixed
polynomial $a$,
\begin{equation}\label{eq:conicseedorder}
 L_i(a)\in e^{kt}r^2\C[[r]].
\end{equation}
This bound holds for all original markers, not just for the unit.

Suppose first $d_F\ge2$. The other critical points of $P$ are
$x=y=0$, $F'(w_p)=0$. Each $w_p$ is nonzero because $h\ne0$.
If it has multiplicity $\mu_p$, analytic coordinates give
\[
 P=k+F(w_p)+s z+y_0^{\mu_p+1}.
\]
Indeed absorb the nonvanishing factor $w$ into $x$ or $y$, then
use the one-variable power coordinate for $F$ at $w_p$.
The explicit characters \eqref{eq:ordinarycharacters}, including
their integration-by-parts and time recurrences, apply to this
exact local model. Their prefactors are
\[
 e^{(k+F(w_p))t}t^{-1-j/(\mu_p+1)},
                    \qquad 1\le j\le\mu_p.
\]
The Gaussian and coordinate Jacobian are analytic units in that
construction. These characters retain multiple roots and equal
critical values. Their total number is $\sum_p\mu_p=d_F-1$.

The $d_F+3$ characters just written are independent on polynomial
amplitudes. Put $B_0(w)=\prod_p(w-w_p)^{\mu_p}$. For the four
conic rows use $x^jB_0(w)$, $0\le j\le3$. Their leading
conic matrix is $B_0(0)I_4$. Their leading finite jets vanish.
For the remaining rows choose polynomials in $w$ with value zero
at $w=0$ and the separate Hermite jets of orders
$0,\ldots,\mu_p-1$ at the nonzero $w_p$. The Chinese remainder
theorem supplies them. Their conic values have an additional
positive power of $r$. Their normalized finite evaluation matrix
is the invertible local jet map, multiplied by analytic units.
After dividing each character column by its scalar prefactor,
the leading matrix is therefore block diagonal and invertible.
Since its size is the proved full rank, these are a full formal
frame. No uncomputed image kernel or infinity block remains.

For $d_F=1$ there is instead one further character. In the exact
coordinates
\[
 x=X/t,\qquad y=Y/t,\qquad w=(h/m)t+W,
\]
the phase is
\begin{align*}
 tP-q={}&\frac{h^2}{2m}t^2+(k-h\ell_w/m)t+H(X,Y,W)+O(t^{-1}),\\
 H={}&\frac hm XY-\frac hm(DX+EY)-\frac m2 W^2-\ell_w W-q_0.
\end{align*}
The correction is polynomial in $t^{-1},X,Y,W$, with positive
$t^{-1}$ order. The Hessian determinant of $H$ is $h^2/m\ne0$.
Formal Gaussian evaluation gives a horizontal character with
nonzero quadratic action and volume prefactor $t^{-2}$.
To check independence set $w_*=(h/m)t$. Use
$x^j(1-w/w_*)$, $0\le j\le3$, and $w/w_*$ as five polynomial
amplitudes with rational parameter coefficients. The former
have the original conic leading matrix and zero leading escape
values; the latter has zero leading conic values and nonzero
leading escape value. This gives the full rank-five frame.

In either degree case, after finite ramification the frame splits
the module into rank-one factors. Every zero-exponential
functional has positive order on every original polynomial
marker, by \eqref{eq:conicseedorder} and the displayed finite
prefactors. Lemma~\ref{lem:positiveformal} rules out a nonzero
constant original marked period. A constant-coefficient operator
reduces any nonzero finite exponential-polynomial to $c e^{\beta t}$;
it replaces the marker by another original polynomial in $P$.
Shifting $k$ by $-\beta$ and normalizing the marker gives the same
contradiction. This proves the theorem.
\end{proof}

\begin{corollary}\label{cor:ordinaryconics}
All the critical-conic and critical-hyperbola rows in the ordinary-node
table are not pure. Proposition~\ref{prop:axisexclusion} treats its remaining center-axis row.
\end{corollary}
\begin{proof}
For $h\ne0$ and $a=b=c=d=e=f=0$ the residual is
$T(uv+h)+k$. Apply the $d_F=1$ case with $w=T$, $x=u$, $y=v$.
The no-frame conditions give $ABm\ne0$.

For $h=0$, $a=b=e=0$, $f\ne0$, permute the coordinates and
write the residual as
\[
 v(Tu+d v^2+c v+f)+k.
\]
Take $w=v$, $x=T$, $y=u$, and
$F(w)=d w^3+c w^2+f w$. The quadratic coefficients along $T,u$
are nonzero: the original $T$ square coefficient is nonzero, and the $u$
node has $F_2=0$, so its Gaussian square coefficient is nonzero
by the no-frame condition. If this $F$ is linear, the remaining
$v$ node also has $F_2=0$, giving the extra $w$ square coefficient
required in degree one. The symmetric hyperbola row is identical
after swapping $u,v$. These are all the indicated rows.
\end{proof}

\subsection{A uniform cut criterion in any degree}
\label{sec:uniformcutdecay}

Intersecting critical curves need not admit a simple formal normal
form. The following criterion avoids choosing one. Its hypothesis
concerns the full phase $tP-q$, uniformly in large complex $t$;
bounded critical points alone would not suffice.

\begin{theorem}\label{thm:uniformcuts}
Let $P,q\in\C[x_1,\ldots,x_n]$, with $P$ nonconstant. Suppose
there are fixed $c,R_0,t_0>0$ such that
\begin{equation}\label{eq:uniformphasegradient}
 \|\nabla_x(tP-q)(x)\|\ge c\|x\|
       \quad(|t|\ge t_0,\ \|x\|\ge R_0).
\end{equation}
Along every fixed ray $t=Re^{i\theta}$, the full rapid-decay
homology has a flat cut basis $\Gamma_i(t)$, attached to the
critical values $c_i(t)$ of $tP-q$, such that for every fixed
polynomial $a$,
\begin{equation}\label{eq:normalizedcutdecay}
 e^{-c_i(t)}\int_{\Gamma_i(t)}a(x)e^{tP-q}\,
                 dx_1\cdots dx_n\longrightarrow0.
\end{equation}
Each $c_i(t)=\alpha_i t+O(1)$, where $\alpha_i$ is a critical
value of $P$. If all critical values of $P$ lie in a closed
half-plane through zero, a supporting $t$ ray makes every flat
period of every such fixed form tend to zero.
\end{theorem}

We give the completeness, tail and limiting-form arguments separately.
None follows just from counting the phase-critical points.

\begin{proof}
For every sufficiently large $t$, \eqref{eq:uniformphasegradient}
makes the phase-critical set finite and contained in one fixed
ball. It also gives complex tameness at infinity. Pham's tame
Milnor decomposition~\cite[Part~I and Appendix]{Pham} identifies
the complete middle rapid-decay homology with the local vanishing
homologies carried along cuts from its critical values. The rank
is the sum of their Milnor numbers, including degenerate points.
If this rank is zero the conclusion is immediate. Otherwise
retain all these local generators; identical critical values
are grouped, not discarded.

On the chosen ray, the critical points and values are algebraic
functions of $R$, after a finite cover. Enlarge $R_1$ so branches
and multiplicities stabilize. Choose $|\eta|=1$ with
$\Re\eta\ge1/2$ so the cuts
\[
 S_t(x)=c_i(t)-\eta s,\qquad s\ge0,\quad S_t=tP-q,
\]
are disjoint for $R>R_1$, except at identical grouped endpoints.
There are only finitely many forbidden choices: every nonzero
difference of two algebraic critical values has one limiting
direction. Choose this auxiliary cut direction after fixing
$\theta$; it imposes no genericity condition on $\theta$ itself.
In particular a supporting $t$ ray need not be a generic ray.

Along a cut, normalized complex gradient transport is
\[
 x_s=-\eta\,\overline{\nabla S_t}/\|\nabla S_t\|^2,
               \qquad S_t(x(s))_s=-\eta.
\]
Outside the fixed ball,
$|d\|x\|^2/ds|\le2/c$. All local vanishing generators start
in that ball. With fixed sufficiently large $A$ and $B>2/c$,
they therefore have representatives inside the carrier
\begin{equation}\label{eq:uniformcarrier}
 \|x\|^2\le A+Bs.
\end{equation}
At its outer boundary the transport points strictly inward
relative to this inequality. No trajectory escapes at finite
cut length, and the cuts avoid all other critical values.

We must control a single family of representatives. The gradient
flow itself need not be semialgebraic. Instead compactify $x$
by $x/\sqrt{1+\|x\|^2}$ and $s$ by $s/(1+s)$. Take the closure
of each constrained cut carrier and mark its rapid boundary
$G$, where $s=\infty$. Inequality~\eqref{eq:uniformcarrier}
forbids spatial infinity at finite $s$. A compatible
semialgebraic triangulation computes the relative homology of
this compact pair by locally finite semialgebraic cut chains,
with boundary allowed only in $G$. These representatives are
admissible rapid-decay chains: on the carrier,
\[
 s\geq(\|x\|^2-A)/B,\qquad
 |e^{-\eta s}|\leq e^{-s/2}
       \leq \exp\!\left(-\frac{\|x\|^2-A}{2B}\right).
\]
Semialgebraic approach to a boundary has at most power-law rates,
so this decay dominates every boundary pole of a polynomial form.
The chain-current boundary identity applies on finite phase bands;
the tail estimate below justifies passage to the improper integral.
Thus the compact-pair calculation supplies admissible cycle
representatives, without identifying arbitrary singular chains
with integration chains. Their exponential is exactly
$e^{c_i-\eta s}$.

For the full unconstrained cut carrier, mark and remove any
spatial boundary at finite $s$; it is not a rapid boundary.
The constrained carrier, the full carrier, their rapid boundaries
and this excluded boundary form a finite family of semialgebraic
subsets in the same compactification. Apply the compatible
finite-family form of Hardt triviality to projection onto $R$.
It gives one trivialization on an unbounded interval, preserving
these subsets and their inclusions. The precise form is stated
in Drusvyatskiy--Lewis~\cite[Theorem~2.25]{DrusvyatskiyLewis};
it is Hardt's theorem. Semialgebraic triangulation and the
chain-current boundary rule are as in
Hardt--Lambrechts--Turchin--Voli\'c~\cite[Sections~2.1 and~3,
especially Theorem~3.5]{HLTV}.

For precision, the pairs in this application are as follows. In the
compactified $(x,s)$ space write $\overline K_i(R)$ for the closure of
the constrained cut carrier, $G_i(R)=\overline K_i(R)\cap\{s=\infty\}$,
and write $\overline K_i^{\rm full}(R)$ for the unconstrained closure.
Let $E_i(R)$ be its spatial boundary at finite $s$, and put
$G_i^{\rm full}=\overline K_i^{\rm full}\cap\{s=\infty\}$. The
relative pairs and their inclusion are
\[
 (\overline K_i,G_i)\longrightarrow
 (\overline K_i^{\rm full}\setminus E_i,
       G_i^{\rm full}\setminus E_i).
\]
The constrained pair has no $E_i$, by \eqref{eq:uniformcarrier}.
At $s=0$ a local vanishing cycle closes by its collapsing thimble;
that endpoint is included in $K_i$, rather than declared a rapid
boundary. In degree $n$, projection to $x$ takes these relative
classes to the rapid-decay homology of $e^{S_t}$. The tame cut
comparison above says that the images, summed over the cuts, contain
the entire vanishing-homology decomposition. We use this image;
we need not identify every relative group of an individual carrier
with the full rapid-decay group.

The compatible Hardt trivialization preserves both pairs, their
inclusion, and hence the induced maps on their relative homology.
Gradient transport already showed that the constrained inclusion
contains every local vanishing generator. At one parameter choose
finitely many integral relative chains whose images are a basis
over $\C$ of the full rapid-decay space. A nonprimitive integral
image causes no problem over $\C$. Their number is the total tame Milnor rank, and their images are
independent at this parameter by construction. Transport these chains by
the compatible trivialization. They form a jointly semialgebraic
family with fixed integer multiplicities and supports satisfying
\eqref{eq:uniformcarrier}. Their images are a flat basis of the
actual homology local system: on each compact parameter interval
the transport preserves the relative pairs and is proper on
every finite phase band. The induced comparison maps commute with this transport, so the image
basis remains independent and spanning. This establishes flatness and
completeness before any mass estimate is made; a uniform bound for
arbitrary representatives would not establish either property.

Here is the needed uniform mass bound. In a fixed semialgebraic
family of $d$-chains with bounded multiplicities, the mass inside
a ball of radius $L\ge1$ is at most $C L^d$, uniformly in the
parameters. Partition tangent planes by their largest coordinate
$d$-minor. The Cauchy--Binet identity gives a lower bound
$\binom Nd^{-1/2}$ for some such minor in ambient dimension $N$.
The area formula then bounds that piece by its coordinate
projection volume times the number of preimages. Semialgebraic
uniform finiteness bounds those finite fibers independently of
the parameters; see, for example,
\cite[Theorem~2.32]{DrusvyatskiyLewis}. Sum the finitely many
pieces and multiplicities. Apply the same argument to boundary
chains. This proves uniformity, rather than merely pointwise
finite volume.

Lift the cut chains to $(x,s)$. Their bands $j\le s\le j+1$
lie in balls of radius $O(1+j)$ by
\eqref{eq:uniformcarrier}. Their fixed semialgebraic descriptions
and the mass estimate give, with constants independent of $R$,
\begin{equation}\label{eq:uniformcuttail}
 \left|e^{-c_i(t)}\int_{s\ge L}a e^{S_t}\,dx_1\cdots dx_n\right|
 \le C\sum_{j\ge\lfloor L\rfloor}(1+j)^M e^{-j/2}.
\end{equation}
The polynomial amplitude changes only $M$. This tail tends to
zero uniformly in $R$.

It remains to bound a fixed phase band. A bounded critical branch
$x(r)$, $r=1/t$, satisfies
\[
 q(x(r))=-c_i(r)-r c_i'(r).
\]
At a stationary point the chain rule gives $c_i'(r)=-P(x(r))/r^2$;
eliminating $P$ gives this identity. Since $q(x(r))$ is
bounded, its Puiseux expansion forces the only possible polar
term of $c_i$ to be $\alpha_i/r$. Its coefficient is
$\alpha_i=P(x(0))$, and $\nabla P(x(0))=0$. Hence
$c_i(t)=\alpha_i t+O(1)$ even at critical collisions.

Choose a generic fixed truncation level $L$. The truncated lifted
chains have compact supports and uniformly bounded mass and
boundary mass. The latter follows from the same semialgebraic
estimate applied to their slices; after a finite stratification
there is an arbitrarily large sequence of such levels. By the
Federer--Fleming compactness theorem, every sequence $R\to\infty$
has a subsequence converging to an integral $n$-current. We use
the version with both mass and boundary-mass bounds,
Simon~\cite[Chapter~6, Section~3, Theorem~3.15]{SimonGMT}.
Its rectifiability conclusion is essential.

On the lifted chain,
\[
 P(x)-\alpha_i=
       \frac{c_i(t)-\alpha_i t+q(x)-\eta s}{t}.
\]
Every limiting current is therefore supported in
$\{P=\alpha_i\}\times[0,L]$. Stratify this complex algebraic
level set into smooth complex strata. Almost everywhere for the
rectifiable limiting current, its tangent plane projects into
the tangent of one of these strata. Each has complex dimension
less than $n$. The holomorphic $n$-form $dx_1\cdots dx_n$ thus
vanishes on that tangent plane. Lower-dimensional strata carry
no $n$-mass; the same argument covers singular strata of larger
real dimension. Testing with the fixed smooth form
$e^{-\eta s}a(x)dx_1\cdots dx_n$, with a compact cutoff, gives
zero on every subsequential limit. Consequently the normalized
integral on $0\le s\le L$ tends to zero. This does not assert
that the real volume of the chains tends to zero.

First let $R\to\infty$, then let $L\to\infty$ through the
allowed levels in \eqref{eq:uniformcuttail}. This proves
\eqref{eq:normalizedcutdecay}. Finally choose $\theta$ with
$\Re(e^{i\theta}\alpha)\le0$ for all critical values $\alpha$
of $P$. The factors $e^{c_i(t)}$ are then bounded, so every basis
period tends to zero. A flat cycle is a constant linear
combination of this full basis, proving the last assertion.
\end{proof}

\begin{corollary}\label{cor:twocriticalvalues}
Let $q$ be a full affine Gaussian quadratic and assume
\eqref{eq:uniformphasegradient}. If $P$ has at most two critical
values, every original marked Gaussian period is either zero or
is not a finite exponential-polynomial. In particular $P$ is
neither pure nor safe.
\end{corollary}
\begin{proof}
Any two complex numbers lie in a closed half-plane through zero.
For two nonzero values normalize them to $z_1,z_2$ on the unit
circle. If $z_1+z_2\ne0$, use the direction
$-\overline{z_1+z_2}/|z_1+z_2|$; if their sum is zero, use
$i\overline{z_1}$. Zero or missing values add no condition.

A nonzero constant original marked period would, by the rational
seed restriction and actual-period extension in
Lemma~\ref{lem:seedfunctional}, be represented by a flat period
of that fixed polynomial form. Its generic exponential system
has $f=t$, $g=tP-q$ on smooth affine space and regular input
$\mathcal O$; $f,g$ are algebraically independent because $P$
is nonconstant. Remove the finite exceptional parameter set of
the Hien--Roucairol comparison. Theorem~\ref{thm:uniformcuts}
makes the resulting constant period tend to zero, a contradiction.
No convergence of the formal moment series was assumed.

A constant-coefficient operator reduces a nonzero finite
exponential-polynomial to $c e^{\beta t}$ while replacing its
original marker by another fixed polynomial in $P$ times that
marker. Shift $P$ by $-\beta$ and normalize. Its gradient is
unchanged and its critical values are still at most two. The
same contradiction applies. Since $\Phi_1(0)=1$, it cannot be
a polynomial, proving also the safety exclusion.
\end{proof}

\subsection{Completing the ordinary-node center axes}
\label{sec:ordinaryaxes}

The only remaining ordinary-node row is
\begin{equation}\label{eq:ordinaryaxes}
 P=Tuv+a u^3+d v^3+b u^2+c v^2+k.
\end{equation}
Keep the full $q$ of Lemma~\ref{lem:ordinarynormal}, with
$m\ne0$. Under the no-frame assumption the other nodes satisfy
\[
 a=b=0\ \Longrightarrow\ A\ne0,\qquad
 d=c=0\ \Longrightarrow\ B\ne0.
\]
We use a full formal frame when $bc\ne0$ and the preceding
uniform cut criterion when $bc=0$. These branches are disjoint
and include every zero face of the four coefficients.

\begin{lemma}\label{lem:axisuniformgradient}
For \eqref{eq:ordinaryaxes} with these Gaussian conditions,
\eqref{eq:uniformphasegradient} holds.
\end{lemma}
\begin{proof}
Otherwise take $\|x\|\to\infty$, $|t|\ge t_0$ and
$\nabla(tP-q)=o(\|x\|)$. Division by $t\|x\|^2$ shows that
a limiting projective direction is a leading node: the $T$ axis,
also the $u$ axis if $a=0$, and the $v$ axis if $d=0$.

Near any such node use height $z$ and transverse coordinates
$x,y=o(z)$. The phase polynomial has the form
\[
 P=zxy+\beta z^2+H(x,y)+k,
\]
where $H$ has only terms of degrees two and three. There are
no transverse linear terms. The transverse equations, divided
by $tz$, give
\[
 |x|+|y|\le C/|t|+o(|x|+|y|)+o(1/|t|).
\]
The terms of $\nabla H/z$ are absorbed because $(x,y)=o(z)$;
the Gaussian cross terms divided by $tz$ are bounded by the
same quantities. Thus $x,y=O(1/|t|)$. The height equation is
\[
 \partial_z(tP-q)=(2\beta t-M)z+O(1),
\]
where $M$ is the Gaussian square coefficient along this node.
If $\beta\ne0$, $|2\beta t-M|$ is uniformly bounded below
for all sufficiently large $|t|$. If $\beta=0$, the no-frame
condition gives $M\ne0$. This contradicts the assumed small
gradient ratio. The proof also covers bounded-$t$ subsequences
within the chosen large-$|t|$ domain and therefore proves the
required uniform estimate.
\end{proof}

\begin{lemma}\label{lem:axisrank}
The full generic twisted rank of \eqref{eq:ordinaryaxes} is
\[
 N=5+\mathbf1_{a\ne0}+\mathbf1_{d\ne0}.
\]
\end{lemma}
\begin{proof}
The three homogenized quadratic gradient equations have total
intersection length eight in $\mathbb P^3$. Their points at
infinity are exactly the leading ordinary nodes, in number
$3-\mathbf1_{a\ne0}-\mathbf1_{d\ne0}$. Each has length one.
Indeed its local Jacobian has the two transverse columns of
$\operatorname{Hess}F_3(w)$ and the column $\nabla(tF_2-q_2)(w)$.
The first two have rank two and are annihilated by pairing with
$w$; the last has nonzero pairing
$2(tF_2-q_2)(w)$. This value is nonzero for large generic $t$
by the stated conditions. Thus these infinity intersections are
simple. The uniform gradient estimate gives a finite affine
critical set; there are no other components. B\'ezout and
Pham's tame comparison give the asserted full rank.
\end{proof}

\begin{proposition}\label{prop:axisexclusion}
Every phase \eqref{eq:ordinaryaxes} under the stated conditions
is not pure. Every original marked Gaussian period is zero or
is not a finite exponential-polynomial.
\end{proposition}
\begin{proof}
First let $bc\ne0$. Set $r=t^{-1/2}$, $u=rU$, $v=rV$, leaving
$T$ unchanged. Exactly,
\begin{align*}
 tP-q={}&k/r^2+F_0(U,V,T)+E_r(U,V,T),\\
 F_0={}&TUV+bU^2+cV^2-mT^2/2-\ell_TT-q_0,\\
 E_r={}&r(aU^3+dV^3-DUT-EVT-\ell_uU-\ell_vV)\\
       &-r^2(AU^2+BV^2+2CUV)/2,\\
 du\,dv\,dT={}&r^2dU\,dV\,dT.
\end{align*}
The full polynomial twisted quotient for $F_0$ has rank five,
including its degenerate parameter values. Here is a separate rank
proof. Suppose $\nabla F_0=o(\|(U,V,T)\|)$ on an unbounded
sequence. If $T$ stays bounded, the first two equations force
both normalized $U,V$ to be nonzero; the last, divided by the
squared norm, gives their product zero, a contradiction. If
$T$ is unbounded, the first two equations have inverse
$O(1/|T|)$, giving $U,V=o(1)$ and $\|(U,V,T)\|\asymp|T|$.
The last equation then has leading term $-mT$, again a
contradiction. This coercivity is uniform on compact sets with
$bcm\ne0$. The coefficient domain is connected, so conservation
of critical intersection multiplicity and Pham's tame comparison
make its rank constant. At generic $\ell_T$, its critical
equations give one point with $U=V=0$ and four with
$T^2=4bc$, all simple. Hence the rank is five, also when points
merge at the exceptional values of $\ell_T$.

Choose any polynomial basis $B_1,\ldots,B_5$ of this fixed
quotient and its dual $\lambda_1,\ldots,\lambda_5$. Then
\[
 L_j(a_0)=e^{kt}r^2\lambda_j\bigl(e^{E_r}a_0(rU,rV,T)\bigr)
\]
defines five formal horizontal functionals on the original full
module. Coefficients are finite polynomial evaluations at every
order. Spatial exactness is the fixed-phase conjugacy; time
horizontality retains the displayed $r^2$ volume and the
coordinate velocity $(u/r,v/r,0)$, whose divergence is $2/r$.
Cartan's formula makes the remaining velocity term spatially exact.
For every original fixed polynomial marker their values lie in
$e^{kt}r^2\C[[r]]$.

The other $P$-critical points are the off-axis points
$p_u=(-2b/(3a),0,0)$ if $a\ne0$ and
$p_v=(0,-2c/(3d),0)$ if $d\ne0$. They are Morse, with Hessian
determinants $8b^3/(9a^2)$ and $8c^3/(9d^2)$. Holomorphic Morse coordinates reduce their local phases to a
nondegenerate quadratic. Expanding the analytic Gaussian amplitude
and volume unit and applying the quadratic Wick recurrence gives
all-order horizontal characters. These have prefactors $e^{P(p)t}r^3$, multiplied by power series in
$r^2$ on fixed markers. Their values may coincide with each
other or with $k$; no separation hypothesis is used.

Here is an explicit completeness test. Let $M$ bound the degree
of the chosen $B_j$ in $U,V$. Substitute $U=u/r,V=v/r$ and
multiply each $B_j$ by
\[
 C(u,v)=\prod_{p_u}(1-u/u_p)^K
                      \prod_{p_v}(1-v/v_p)^K,\qquad K>M.
\]
An empty product is one. At the central fixed model this is
$1+O(r)$, so the normalized five-by-five matrix is the identity
at leading order. At each off-axis point its zero of order $K$
overcomes the worst $r^{-M}$ growth of the substituted basis;
the normalized cross values tend to zero. For the off-axis rows
use $(u/u_p)^K$ or $(v/v_p)^K$. They have zero central leading
values, unit value at their own Morse point and zero leading
value at the other. The resulting full matrix has invertible
block diagonal leading term. Its size is exactly the rank of
Lemma~\ref{lem:axisrank}. These are therefore a full formal
frame, including all fixed-model degeneracies and equal actions.

Every zero-exponential original-marker seed has positive order,
either $r^2$ or $r^3$. Lemma~\ref{lem:positiveformal} excludes
a nonzero constant original marked period. Differential
isolation and shifting $k$ then exclude every nonzero finite
exponential-polynomial as before.

It remains to treat $bc=0$. The critical equations of $P$ show
that its values are among
\[
 k,\qquad k+4b^3/(27a^2)\quad(ab\ne0),\qquad
 k+4c^3/(27d^2)\quad(dc\ne0).
\]
At a missing isolated point omit its displayed value. Any extra
critical axes have value $k$. Since $bc=0$, there are at most
two values, including every intersecting-axis or degenerate
case. Lemma~\ref{lem:axisuniformgradient} and
Corollary~\ref{cor:twocriticalvalues} apply. This completes both
disjoint branches and proves the proposition.
\end{proof}

\begin{corollary}\label{cor:ordinarycomplete}
Frame extraction is proved for all three ordinary-node leading
types: irreducible nodal cubics, transverse line--conic unions,
and three nonconcurrent lines. In fact none of these leading
types permits a pure polynomial.
\end{corollary}
\begin{proof}
If all leading nodes have $F_2\ne0$, use the coercive-node
criterion. Otherwise an isotropic zero-$F_2$ node supplies the
frame directly. In the absence of such a frame, the common
normal form applies and satisfies the no-frame conditions at
every other node. The exhaustive critical-set table is now
entirely covered by Proposition~\ref{prop:ordinaryfinite},
Corollary~\ref{cor:ordinaryconics}, and
Proposition~\ref{prop:axisexclusion}. Thus a pure polynomial
must have a frame.

In that frame, the quadratic coefficient of $W$ in the cubic
leading form is the tangent quadratic at the chosen ordinary
node. It has rank two, so its discriminant is nonzero.
Proposition~\ref{prop:nondeglog} excludes purity in that branch
as well. This proves the stronger assertion.
\end{proof}

\subsection{It is enough to complete the regular part}
\label{sec:boundedregular}

The cusp and tangent cases share a substantial simplification.
We need not construct every escaping Gaussian character. It is
enough to exclude escaping contributions to the regular part,
compute that part's full rank, and construct precisely that many
local functionals. This keeps the original marker while avoiding
an unnecessary classification of the other exponential blocks.

\begin{lemma}[A positive seed on the regular part]
\label{lem:regularpositiveseed}
Let $M$ and a seed $a$ satisfy the formal Gaussian hypotheses of
Lemma~\ref{lem:seedfunctional}. After a finite cover at infinity,
write $M_0$ for the zero-exponential summand. Evaluate its dual
functionals on the canonical projection of $a$ to $M_0$. Suppose $M_0$ has a
full horizontal dual frame over $\C((\rho))$, and every member
takes $a$ into $\rho\C[[\rho]]$. Then the formal Gaussian
functional cannot take $a$ to one. No semisimplicity assumption
on the nonzero-exponential summands is required.

The same conclusion holds if the regular frame exists over
$\C((\rho))(\log\rho)$ and its seed values are finite
polynomials in $\log\rho$ with coefficients of strictly positive
$\rho$ order.
\end{lemma}
\begin{proof}
The rational seed restriction supplies a quotient $U\to\C(t)$.
The formal exponential decomposition is functorial and exact on
submodules and quotients; different exponential types have no
horizontal morphisms. These are the classical formal decomposition
facts, recalled in Sabbah~\cite[Section~3, Corollary~3.3 and its
proof]{SabbahStationary}. Consequently the quotient can only use
$U_0\subset M_0$. Since $M_0$ is trivial over the displayed field,
it is semisimple there, and this quotient extends to a horizontal
functional on $M_0$. It is a constant linear combination of its
full dual frame. Positive seed orders cannot sum to one.

Adjoining $\log\rho$, with derivative $1/\rho$, does not add
constants or negative-power exponential solutions. The same
argument applies over this differential field. A finite sum of
positive powers times logarithmic polynomials still cannot be
one. This also explains why a full regular frame, rather than
semisimplicity of the whole ambient module, is the relevant
hypothesis.
For an elementary check, expand this field in
$\C(\log\rho)((\rho))$. The operator $\rho\partial_\rho$
preserves its valuation. A nonzero exponential twist has a
negative-valuation scalar term in that operator, which cannot
be cancelled by a regular matrix. Its horizontal maps to the
trivial module therefore vanish even over this field. Its
constants are $\C$, since a rational function $a(z)$ satisfying
$a'+na=0$ is zero for integer $n\ne0$ and constant for $n=0$.
\end{proof}

\begin{lemma}[No finite-target contribution from spatial infinity]
\label{lem:nospatialinfinity}
Fix a small nonzero parameter $\epsilon_0$. Suppose in a neighbourhood
of $\epsilon_0$ that $\|\nabla_x S\|\ge c\|x\|$ outside a compact set,
where $S=P/\epsilon-q$ and $q$ is quadratic. The outer fibre pairs of
$(\epsilon,S)$ are locally trivial over bounded target neighbourhoods.
Consequently no additional finite-target discriminant or conormal
multiplicity arises from spatial infinity there.
\end{lemma}
\begin{proof}
On $S=g$ with $g$ bounded,
$\partial_\epsilon S=-(g+q)/\epsilon=O(1+\|x\|^2)$,
locally uniformly in $\epsilon$. A target-lifting vector field obtained
by dividing the conjugate gradient by its squared norm has size
$O(1/\|x\|)$. Correcting a parameter-lifting vector field by this
field has size $O(1+\|x\|)$. After a cutoff inside a larger ball,
these fields are complete for finite transport times: the radial
growth inequality is at most linear. They lift the coordinate paths
in the target and parameter neighbourhood and identify the exterior
fibre pairs. All changes at a finite target are therefore supported
inside a compact set. This is a fixed nonzero-parameter statement;
it does not assert a uniform compact set as $\epsilon\to0$.
\end{proof}

\begin{proposition}[Bounded critical branches determine the regular rank]
\label{prop:boundedregularrank}
Let $P$ have degree $D\ge2$ and leading form $F_D$, and let $q$
be quadratic. Suppose, locally uniformly for small nonzero
$\epsilon$, the phase $S=P/\epsilon-q$ has a complex linear
gradient lower bound outside a compact set. The constants and
compact set may depend on $\epsilon$. Assume also
\begin{equation}\label{eq:nonisoleading}
 q_2(w)\ne0\quad\hbox{whenever }[w]\in\mathbb P^{n-1}
                           \hbox{ and }\nabla F_D(w)=0.
\end{equation}
Then the rank of the full zero-exponential summand equals the sum
of the generic local Milnor numbers of bounded phase-critical
branches whose limiting $P$ value is zero. An escaping branch
never has action polar part $\alpha/\epsilon$ for any constant
$\alpha$.
\end{proposition}
\begin{proof}
At a stationary branch the chain rule gives
$S'=-P/\epsilon^2$, hence $q=-S-\epsilon S'$.
A bounded branch has bounded $q$, so its action has only the
possible polar term $P(x(0))/\epsilon$. If a branch escapes,
divide $\nabla P=\epsilon\nabla q$ by $\|x\|^{D-1}$.
Its limiting direction is singular for $F_D$.
Condition~\eqref{eq:nonisoleading} makes the quadratic term of
$q$ along that direction nonzero. Thus $q(x(\epsilon))$ has a
pole, whereas $-S-\epsilon S'$ would be bounded for either zero
polar part or a single term $\alpha/\epsilon$. This excludes
both possibilities, without an escape-scale or Morse assumption.

We spell out the global rank comparison. Fixed-parameter tameness
gives full twisted rank equal to the total critical Milnor number
by Pham's decomposition. For the regular algebraic input
$\mathcal O$ on $\C^*_{\epsilon}\times\C^n$, form the ordinary
direct image by $(\epsilon,S)$. Its relevant cohomology is the
middle reduced Gauss--Manin module; the constant fiber cohomology
has zero exponential integral in the affine $S$ coordinate.
Lemma~\ref{lem:nospatialinfinity} excludes any finite-target
contribution from spatial infinity, using the locally uniform gradient
hypothesis. Every nonvertical discriminant curve at finite target comes
from an affine critical value, with conormal multiplicity equal
to the sum of its local Milnor numbers.

First ramify all critical-value curves into graphs. Roucairol's
nonramified rank formula, Theorem~0.1(2), Section~1.2, together
with Theorem~0.3~\cite{RoucairolFormal}, computes every nonzero
exponential-block rank from those graph multiplicities. Count
all conjugate graphs even when their polar parts agree. No
distinct-action assumption or monodromy formula is used.
More explicitly, on the ramified parameter disc let $\ell$ index
the distinct critical-value graphs $c_\ell(\epsilon)$ and put
\[
 \mu_{\rm tot}=\sum_x\mu_x(S),\qquad
 m_\ell=\sum_{x:\,S(x)=c_\ell}\mu_x(S),\qquad
 Q_\ell=\operatorname{pp}_{\epsilon=0}c_\ell.
\]
The first sum is over the generic affine critical scheme, and the
second over the branches on that value graph. Tameness gives
$\operatorname{rank}M=\mu_{\rm tot}=\sum_\ell m_\ell$.
The conormal multiplicities entering Roucairol are exactly $m_\ell$:
the affine Milnor decomposition gives them locally and the preceding
lemma excludes an extra term at spatial infinity. In the graph case
the rank formula reads
\[
 \operatorname{rank}M_Q=\sum_{\ell:Q_\ell=Q}m_\ell\quad(Q\ne0),
 \qquad
 \operatorname{rank}M_0
 =\mu_{\rm tot}-\sum_{Q\ne0}\operatorname{rank}M_Q
 =\sum_{\ell:Q_\ell=0}m_\ell.
\]
All conjugate graphs are included after the cover; equal polar parts
are grouped only in the sum for $M_Q$. Exactly the bounded-action
branches remain. By the first part
they are precisely the bounded branches with limiting $P$ value
zero. This proves the assertion about the entire regular rank.
\end{proof}

We use the following explicit way to construct that regular part.
Suppose a bounded group at $P$ value $\alpha$ has coordinates,
after a cover, in which
\begin{equation}\label{eq:boundedfixedmodel}
 S=\alpha t+H_0(z)+E_\rho(z),\qquad
 \omega=\rho^\kappa J_\rho(z)\,dz,
 \quad E_\rho\in\rho\C[z][[\rho]],\quad\kappa>0,
\end{equation}
where $J_\rho$ is a unit with polynomial coefficients at each
order. If $\lambda$ is a dual functional on the fixed twisted
quotient of $H_0$, then
\[
 L(A)=e^{\alpha t}\rho^\kappa\lambda\!\left(
           e^{E_\rho}J_\rho A(x(\rho,z))\right)
\]
is a formal functional on the original module. The product rule
proves spatial exactness. Cartan's formula for the coordinate
velocity, together with the displayed moving Jacobian, proves
parameter horizontality. Every coefficient uses finitely many
polynomial terms. Fixed original markers with nonnegative powers
in these coordinates have positive seed order.

The rank check has two directions. First, every bounded critical
branch in the group must have bounded scaled coordinates $z$.
Its limiting $z$ then belongs to the critical scheme of $H_0$,
since $d_zE_\rho$ has positive order. Conversely, every isolated
critical point of $H_0$ deforms with its full local intersection
multiplicity in $H_0+E_\rho$. In the applications below the
coordinate map $x(\rho,z)$ is regular at $\rho=0$ near these
points, so all the deformed points give bounded original
critical branches. No root of the fixed model is merely counted
without being realized in the original family. Conservation of
local intersection multiplicity therefore counts exactly the
whole group, including collisions. One must also check that the
fixed critical scheme is finite and that its total Milnor number
equals the fixed twisted rank.

A polynomial basis of the fixed quotient then gives independent
functionals at leading order. Distinct groups are separated by
polynomial jets, with arbitrarily high vanishing orders to overcome
the finite negative powers introduced by the scaled representatives.
When their number equals the regular rank of
Proposition~\ref{prop:boundedregularrank}, they exhaust $M_0$.
These checks will be made explicitly below; the formal perturbation
formula alone does not establish completeness.

\subsection{The cusp and tangent families together}
\label{sec:cusptangent}

Both remaining irreducible or line--conic cases have a unique
leading singular direction. If $F_2$ is nonzero there, a gradient
argument excludes purity. If it is zero and the direction is
isotropic, the frame test applies. We treat the nonisotropic
alternative with its full Gaussian
\[
 q=\tfrac12(Au^2+Bv^2+mT^2+2Cuv+2DuT+2EvT)
                  +\ell_u u+\ell_v v+\ell_TT+q_0,
 \qquad m\ne0.
\]
Translations and linear changes give the two exact normal forms
\begin{align}
 P_C&=T(u^2+\gamma v+h)+v^3+buv+eu+fv+k,
                       &&\hbox{cusp},\label{eq:cuspnormal}\\
 P_T&=T(u^2+gv+h)+(u+c)v^2+eu+fv+k,
                       &&\hbox{tangent line--conic}.\label{eq:tangentnormal}
\end{align}
Here is the leading-form reduction. At the unique singular
point, $F_3$ is $T$ times a rank-one binary quadratic plus a
binary cubic. Choose that quadratic as $u^2$. A linear shift
of $T$ removes the $u^3,u^2v$ terms, leaving
$Tu^2+a uv^2+b_0v^3$. If $b_0\ne0$, a shift of $v$ by a
multiple of $u$ removes $uv^2$, and rescaling gives the cusp.
If $b_0=0$, tangency to a smooth conic requires $a\ne0$;
rescaling gives the tangent form. The direction of $T$ is
unchanged, so its Gaussian square coefficient remains nonzero.
For the cusp, translate $u$ to remove $uT$, translate $T$ to
remove $u^2$, and translate $v$ to remove $v^2$. For the tangent
case use the same first two translations, then a $v$ translation
to remove $uv$. These operations leave the displayed arbitrary
coefficients and transform, rather than reset, the means and
cross terms of $q$.

\begin{lemma}\label{lem:cusptangentgate}
For both leading forms, $F_2$ nonzero at the singular direction
implies $\|\nabla P\|\ge c\|x\|$ outside a complex compact set.
In the nonisotropic forms above, each fixed nonzero phase
$tP-q$ has the same type of lower bound, locally uniformly in
$t$. Their full generic twisted ranks are six and five,
respectively.
\end{lemma}
\begin{proof}
Any failure of a linear gradient bound approaches the unique
leading singular direction, so $u,v=o(T)$. For the cusp,
the $u$ equation divided by $T$ makes $u$ bounded. The $T$
equation divided by $T$ then has the nonzero limit $2F_2(0,0,1)$.
For the tangent, the transverse equations give
$uv=O(T)$ and $2Tu+v^2=O(T)$. If $u$ is unbounded these imply
$u=O(|T|^{1/3})$; otherwise $u^2/T\to0$ directly. The $T$
equation again has the same nonzero limit. Apply the Gaussian
gradient criterion. Replacing the quadratic value by $-m/2$
gives the fixed-phase bounds for the normal forms. Compactness
of the coefficient estimates gives local parameter uniformity.

For clarity, the ranks do not come from an escape count. The
three quadratic gradient equations have total projective length
eight. Their only infinity point is the $T$ direction. In the
chart $T=1$ with homogenizing coordinate $z$, the cusp $u$
equation solves $u=O(z)$; the $T$ equation then gives $z=0$,
and the last equation is $3tv^2=0$. Its length is two. For
the tangent the $T$ equation solves $z=O(u^2)$, the $u$ equation
gives $u=-v^2/2+O(v^4)$, and the $v$ equation has leading term
$-tv^3$. Its length is three. Fixed-phase coercivity excludes
affine components of positive dimension. B\'ezout and Pham's
tame comparison therefore give ranks $8-2=6$ and $8-3=5$.
\end{proof}

Every escape in these two normal forms approaches the $T$
direction, where $q_2=m/2\ne0$. Proposition~\ref{prop:boundedregularrank}
therefore applies. We now list all possible bounded contributions;
the other exponential blocks require no explicit construction.

\subsubsection*{Isolated critical points on a smooth constraint}

Write $P=T A(u,v)+G(u,v)$. Where $A=0$ is smooth, take $s=A$
and a coordinate $y$ along that curve. The analytic shear
\[
 T'=T+\frac{G(y,s)-G(y,0)}s
\]
gives the exact identity $P=sT'+G(y,0)$. At an isolated critical
point this becomes $\alpha+sT'+y^{\mu+1}$. The already explicit
characters \eqref{eq:ordinarycharacters} provide all $\mu$
dimensions with positive original-marker order.

For the cusp, $\gamma\ne0$ gives a degree-six restriction to
$v=-(u^2+h)/\gamma$, so total isolated multiplicity is five.
If $\gamma=0,h\ne0$, the two lines $u=\pm\sqrt{-h}$ each
have a cubic restriction in $v$, giving multiplicity two each.
For the tangent, $g\ne0$ gives a degree-five restriction and
multiplicity four. If $g=0,h\ne0$, each line $u=u_i$ gives
\[
 (u_i+c)v^2+fv+eu_i+k.
\]
It has one Morse point when $u_i+c\ne0$, none when
$u_i+c=0,f\ne0$, and a critical parabola when both vanish.
The last case is included below. All repeated roots and equal
critical values in the isolated cases retain their Milnor
multiplicities; no Morse perturbation is assumed.

\subsubsection*{Four fixed models along the $T$ lines}

Here and below a scalar constant is harmless but retained in the
phase value. Set $\epsilon=1/t$. The critical equation in $T$
is
\begin{equation}\label{eq:cuspcriticalbound}
 u^2=\epsilon(mT+Du+Ev+\ell_T)
\end{equation}
whenever the constraint is the double line $u^2=0$.
Thus every bounded phase-critical branch has $u=O(\epsilon^{1/2})$.
This estimate will also verify that the following scalings lose
no bounded contribution.

First consider a critical $T$ line at $u=0,v=v_0$ whose transverse
quadratic part is $Tu^2+\lambda(v-v_0)^2+\beta u(v-v_0)$,
with $\lambda\ne0$. Scale $u=t^{-1/2}U$,
$v=v_0+t^{-1/2}V$, leaving $T$ fixed. The volume is $t^{-1}$
and the fixed phase is
\[
 H_0=TU^2+\lambda V^2+\beta UV-mT^2/2-\eta T-q_*.
\]
Complete $V$, then $T$. This is a quartic in $U$, with leading
coefficient $1/(2m)$, plus two nondegenerate Gaussians. Its
twisted rank and critical Milnor number are both three, including
all quartic collisions. Equation~\eqref{eq:cuspcriticalbound}
and the $v$ equation with nonzero linear coefficient $2\lambda$
give $v-v_0=O(t^{-1/2})$. Every bounded branch in this group has
bounded scaled variables. The model therefore supplies all three
local dimensions, with original-marker values in
$e^{\alpha t}t^{-1}\C[[t^{-1/2}]]$.

If instead $\lambda=0,\beta\ne0$, the same scaling gives
\[
 H_0=TU^2+\beta UV-mT^2/2-\eta T-q_*.
\]
The polynomial shear $V\mapsto V+TU/\beta$ removes $TU^2$.
This is a nondegenerate three-variable Gaussian, of rank one.
The $u$ equation with coefficient $\beta$ bounds $v-v_0$ by
$O(u)+O(\epsilon)$, so bounded branches are again captured.
It gives one dimension with the same positive volume order.

For the cusp's remaining line, $P=T u^2+v^3+k$, use
$u=t^{-1/2}U$, $v=t^{-1/3}V$. The fixed phase is
\[
 TU^2+V^3-mT^2/2-\ell_TT-q_0.
\]
Completing $T$ gives a quartic, a cubic and one Gaussian, of
total rank and Milnor number $3\cdot2=6$. The volume is
$t^{-5/6}$. Bounded critical branches satisfy $u=O(t^{-1/2})$
and, from $3v^2=\epsilon q_v$, $v=O(t^{-1/2})$, so their
scaled coordinates are bounded.

For the tangent's remaining line, $P=T u^2+uv^2+k$, use
$u=t^{-1/2}U$, $v=t^{-1/4}V$. Completing $T$ gives the fixed
two-variable phase
\[
 U^4/(2m)+UV^2-\ell_TU^2/m
\]
and one Gaussian. Its rank and critical Milnor number are five.
Here is a direct check: give $U,V$ weights $1/4,3/8$. The
leading derivatives are $2U^3/m+V^2$ and $2UV$, with sole
common zero the origin and quotient basis
$1,U,U^2,V,V^2$. They form a regular sequence. The associated
graded twisted de Rham complex is their Koszul complex,
concentrated in top cohomology; all derivative and lower-phase
terms have smaller weight. It follows that the same five
classes give the twisted rank. The ordinary gradient quotient
has the same weighted leading calculation, so its total
Milnor number is also five for every $\ell_T$. The volume is
$t^{-3/4}$. Equation~\eqref{eq:cuspcriticalbound} bounds $u$,
and $2Tu+v^2=\epsilon q_u$ gives $v=O(t^{-1/4})$. This
captures every bounded branch.

All remaining terms in these four scaled phases have positive
order in a common finite cover. Their coefficients are
polynomials at each order, and every fixed original marker has
nonnegative order. Formula~\eqref{eq:boundedfixedmodel} constructs
the stated complete local functionals. The Jacobian factors and
coordinate-velocity terms are retained, exactly as in its proof.

\subsubsection*{The critical parabola}

This occurs only for the tangent form, with $g=0$, $h\ne0$,
$c^2+h=0$, $f=0$, $c\ne0$. Put
\[
 x=u+c,\qquad s=(x-2c)T+v^2+e,\qquad\alpha=k-ec.
\]
Exactly $P=\alpha+xs$. Write $x=\rho X$, $s=\rho Y$,
$\rho=t^{-1/2}$. The volume is
\[
 \frac{\rho^2}{\rho X-2c}\,dX\,dv\,dY,
\]
and the fixed phase is $XY+\psi(v)$, where
\[
 \psi(v)=-q\!\left(-c,v,\frac{v^2+e}{2c}\right).
\]
Its quartic coefficient is $-m/(8c^2)\ne0$, so its fixed
twisted rank and critical Milnor number are three. At every
bounded phase-critical branch near this parabola, the exact
normal equations give $x=\epsilon q_s=O(\epsilon)$ and
$s=\epsilon q_x=O(\epsilon)$; the denominators stay away
from zero. Thus $X,Y$ are bounded. The scaled Jacobian is a
unit times $\rho^2$, and its expansion and that of the phase
have polynomial coefficients at each order. This gives all
three dimensions, with strictly positive original-marker order.
The other constraint line has its one isolated Morse point.

\subsubsection*{Exhaustive routing and completeness}

The remaining double-line cases require only the following two
tables. Their rows are disjoint. A missing common root means
there is no bounded contribution at all, not an omitted escape.

\begin{center}\small
\begin{tabular}{@{}p{0.36\textwidth}p{0.54\textwidth}@{}}
\toprule
Cusp with $\gamma=h=0$ & Bounded $P$-critical groups\\
\midrule
$f\ne0$ & For each root $3v_i^2+f=0$ with $bv_i+e=0$,
one line with $\lambda=3v_i$, $\beta=b$, rank three\\
$f=0$, $e\ne0$ & None, for every $b$\\
$f=e=0$, $b\ne0$ & One line with $\lambda=0$, $\beta=b$, rank one\\
$f=e=b=0$ & Quartic--cubic model, rank six\\
\bottomrule
\end{tabular}
\end{center}
The line values in the first row are
$\alpha_i=v_i^3+fv_i+k$; the other line values are $k$.

\begin{center}\small
\begin{tabular}{@{}p{0.36\textwidth}p{0.54\textwidth}@{}}
\toprule
Tangent with $g=h=0$ & Bounded $P$-critical groups\\
\midrule
$c\ne0$ & Put $v_0=-f/(2c)$. None if $e+v_0^2\ne0$;
otherwise one rank-three line with $\lambda=c$, $\beta=2v_0$,
value $k-f^2/(4c)$\\
$c=0$, $f\ne0$ & None, for every $e$\\
$c=f=0$, $e\ne0$ & Two lines $v_i^2=-e$, each with
$\lambda=0$, $\beta=2v_i$, rank one, value $k$\\
$c=f=e=0$ & Weighted quartic model, rank five, value $k$\\
\bottomrule
\end{tabular}
\end{center}

To verify these tables directly, the double constraint forces
$u=0$. The cusp equations reduce to $3v^2+f=0$, $bv+e=0$;
the tangent equations reduce to $2cv+f=0$, $v^2+e=0$.
The remaining $T$ coordinate is free. Together with the smooth
constraint and separated-line alternatives already proved,
this exhausts all coefficients of both normal forms.

For every fixed model the critical scheme is finite with the
stated total Milnor number. The bounds above show that every
bounded phase-critical branch has bounded scaled coordinates.
Conservation of local critical multiplicity thus gives exactly
the stated group counts. Proposition~\ref{prop:boundedregularrank}
shows that the groups at value zero account for the entire
regular rank; all escaping branches are excluded from it.

We also need independence on the original polynomial module.
Within a fixed group evaluate a polynomial basis of its fixed
quotient: its normalized leading matrix is invertible by
\eqref{eq:boundedfixedmodel}. Isolated groups use their finite
power jets. Different line groups have distinct constant $v_i$;
the parabola and the other isolated group have distinct constant
$u$ values. High-order Hermite cutoffs in these coordinates
separate the groups. Choose their orders to exceed the largest pole order of the
finitely many basis representatives, measured in the common cover.
Then all cross evaluations have positive order. In a smooth
isolated group the same assertion is ordinary finite jet
interpolation. At the parabola chart, the leading denominator
$-2c$ is a unit; multiplying by a high enough power of
$u-c$ separates the other line without changing that unit.
These constructions give an invertible block leading matrix
of exactly the regular rank. There is no cyclic-seed assumption
and no appeal to a count of unconstructed escape characters.

After a common cover this regular part is trivial, with all
original-marker values of positive order: $t^{-1}$ in the
line and parabola models, $t^{-5/6}$ in the cusp's last model,
$t^{-3/4}$ in the tangent's last model, and the positive powers
of the isolated characters. Lemma~\ref{lem:regularpositiveseed}
excludes a constant original marked period. Differential
isolation and shifting $k$ exclude every nonzero finite
exponential-polynomial in the nonisotropic normal forms.

\begin{proposition}\label{prop:cusptangentcomplete}
A pure cubic with cuspidal leading curve or tangent line--conic
leading curve has an isotropic singular direction with $F_2=0$.
It satisfies the mixed cutoff $m>3\deg Q$.
\end{proposition}
\begin{proof}
The gate excludes nonzero $F_2$ at the singular direction.
The exhaustive bounded-group construction excludes the
nonisotropic alternative, including every collision and
degree drop. In the remaining isotropic alternative,
Lemma~\ref{lem:frame} supplies the Gaussian affine-linear
coordinate. The already proved affine-linear theorem gives
the asserted cutoff.
\end{proof}

\subsection{Three concurrent distinct leading lines}
\label{sec:concurrentregular}

Let the common vertex be the $T$ direction, so the leading form
is a squarefree binary cubic $F(u,v)$. Its binary gradient has
only the zero at the origin. If $F_2$ at the vertex is nonzero,
a sequence with $\nabla P=o(\|x\|)$ must approach that vertex;
then $u,v=o(T)$ and the $T$ equation has a nonzero linear
leading term. The Gaussian gradient criterion excludes purity.
An isotropic zero-$F_2$ vertex already gives a frame.

In the remaining alternative retain the full Gaussian with
$m=q_{TT}\ne0$ and write
\[
 P=T H(u,v)+G(u,v),\qquad\deg H\le1,\quad G_3=F,
 \qquad q=mT^2/2+T J(u,v)+q_*(u,v),
\]
where $J$ is affine linear and $q_*$ includes all lower terms.
For each fixed $t$, the phase gradient is linearly coercive:
its only possible limiting singular direction is the vertex,
where its $T$ derivative has leading term $-mT$. The estimate
is locally uniform on compact parameter sets. Completing $T$
exactly gives
\[
 tP-q=-\frac m2\left(T-\frac{tH-J}{m}\right)^2
                    +tG-q_*+\frac{(tH-J)^2}{2m}.
\]
This change has volume Jacobian one. The binary phase has
squarefree cubic leading form $tF$, whose two leading quadratic
derivatives are a regular sequence of length four. Weighted
leading reduction, or the same Koszul argument used above,
gives full twisted rank four. The Gaussian factor has rank one.
Condition~\eqref{eq:nonisoleading} holds at the sole leading
singular direction. Thus all escapes are excluded from every
pure linear-action block by Proposition~\ref{prop:boundedregularrank}.

If $H$ is a nonzero constant, $P$ has no critical point.
The full regular rank is zero and a constant original marked
period is impossible. This remains true after a constant shift
of $P$.

If $H=0$, integrate the $T$ Gaussian formally. A fixed polynomial
marker becomes a fixed polynomial in $u,v$, using conditional
mean $-J/m$ and covariance $1/m$. The effective Gaussian is
\[
 q_{\rm eff}=q_*-J^2/(2m).
\]
Its Hessian is nondegenerate by the Schur determinant identity
$\det\operatorname{Hess}q=m\det\operatorname{Hess}q_{\rm eff}$.
The squarefree leading binary cubic gives
$\|\nabla G\|\ge c\|(u,v)\|^2$ outside a compact set.
The all-degree Gaussian gradient criterion therefore excludes
every nonzero finite exponential-polynomial original period.
The conditional normalization is independent of $t$.

Finally let $H$ be nonconstant and use $u=H$ as an affine
coordinate. Then $P=uT+G(u,v)$. If $R(v)=G(0,v)$ is
nonconstant, its critical points give all the bounded groups.
The exact smooth-constraint shear gives their local types
$\alpha+s z+y^{\mu+1}$. This includes the empty critical set
when $R$ is linear. Their full power characters have positive
original-marker orders, so the regular-rank and seed lemmas
exclude a constant period without constructing the escaping
blocks.

If $R=k$ is constant, write $G=k+uB(u,v)$, where $B$ has
degree at most two. The critical curve is
\[
 u=0,\qquad T=-B(0,v).
\]
The quadratic coefficient $c$ of $B(0,v)$ is nonzero.
Otherwise the squarefree cubic $F$ would be divisible by $u^2$.
Put $s=T+B(u,v)$, so exactly $P=k+us$ and the coordinate
Jacobian is one. Scale $u=\rho U$, $s=\rho V$,
$\rho=t^{-1/2}$. The fixed phase is
\[
 UV-q(0,v,-B(0,v)),
\]
with quartic coefficient $-mc^2/2\ne0$ in $v$. Its rank and
critical Milnor number are three. The exact critical equations
$u=\epsilon q_s$, $s=\epsilon q_u$ show that every bounded
branch has bounded scaled coordinates. This is precisely the
complete quartic parabola construction above, now with
Jacobian $\rho^2$. It supplies the entire regular block when
$k=0$, with positive original-marker orders; when $k\ne0$
that block has nonzero exponential part. The polynomial
representatives of $v$ powers are original markers, so its
leading matrix is invertible. The seed lemma applies.

These alternatives exhaust the affine polynomial $H$ and the
restriction $R$. They also exclude nonzero finite
exponential-polynomial markers after differential isolation and
shifting the constant term. We obtain the following conclusion.

\begin{proposition}\label{prop:concurrentcomplete}
Every pure cubic whose leading form is three distinct concurrent
lines has an isotropic vertex with $F_2=0$. Hence it has a
Gaussian affine-linear coordinate and satisfies $m>3\deg Q$.
\end{proposition}
\begin{proof}
The gradient gate and the nonisotropic cases just proved leave
only that vertex alternative. Apply Lemma~\ref{lem:frame} and
the complete affine-linear theorem.
\end{proof}

\subsection{Quadratic elimination and the repeated lines}
\label{sec:repeatedlines}

The final two leading types have a repeated factor. They reduce
to a polynomial in one variable after eliminating two Gaussian
variables. The full leading-coefficient theorem makes an escape
table unnecessary: only the actual degree, its leading
coefficient and the determinant factor are needed.

\begin{lemma}[The determinant factor in quadratic elimination]
\label{lem:quadraticelimination}
Over $\C(t)$, let a phase have the form
\[
 -\tfrac12 y^T M(t)y+L(t,x)^Ty+V(t,x),
\]
where $M$ is an invertible symmetric matrix and $L,V$ are
polynomial in $x$. Eliminating $y$ identifies its top twisted
module and constant volume seed with the one-variable phase
and normalization
\[
 \varphi=V+\tfrac12L^TM^{-1}L,\qquad
 h=(\det M)^{-1/2}.
\]
The parameter connection is
$\partial_t+\varphi_t+h'/h$. An algebraic square root is allowed.
If the actual degree $d$ of $\varphi$ is at least two and
$d\,v_\infty(h)>v_\infty([x^d]\varphi)$, its constant seed
cannot have a rational unit quotient.
\end{lemma}
\begin{proof}
The change $y=z+M^{-1}L$ has Jacobian one. Its spatial twisted
relations reduce the $z$ factor to the one-dimensional Gaussian
quotient, with normalized contractions
$\ell(z_i z_j)=(M^{-1})_{ij}$. Parameter differentiation of its
constant class inserts $-z^TM'z/2$, hence
\[
 -\tfrac12\operatorname{tr}(M^{-1}M')=h'/h.
\]
Terms from the moving center are spatial exact forms, by the
product rule or Cartan's formula. This proves the identification
of both connection and seed. After a finite cover at infinity,
$h$ is a Laurent series and Theorem~\ref{thm:stationaryballs}
applies. The cover multiplies both sides of the valuation
inequality by the same positive integer.

The elimination is over the generic field. It need not define
a conditional formal average at $t=0$. The original Wick
functional first gives its rational seed quotient by
Lemma~\ref{lem:seedfunctional}; this quotient is transported
through the displayed module identification. No image-kernel
calculation or convergent original period is needed.
\end{proof}

\begin{lemma}\label{lem:quadraticfibrehighdegree}
Let $u,Y_1,\ldots,Y_s$ be independent standard Gaussian variables
and
\[
 P=Y^TAY+B(u)^TY+g(u),\qquad A=A^T,
\]
with $B,g$ arbitrary polynomials. Put
\[
 D(t)=\det(I-2tA),\quad C(t)=(I-2tA)^{-1},\quad
 \varphi=tg-u^2/2+(t^2/2)B^TCB.
\]
If $\deg_u\varphi\ge3$, then $P$ is not pure.
\end{lemma}
\begin{proof}
Quadratic elimination gives $h=D^{-1/2}$. Let $d=\deg D$ and
$r=\deg_u\varphi\ge3$. Its leading coefficient is $N(t)/D(t)$,
where $N$ is a nonzero polynomial divisible by $t$:
the degree-$r$ coefficient vanishes at $t=0$. Consequently
\[
 r v_\infty(h)-v_\infty(N/D)
     =(r/2-1)d+\deg N>0.
\]
Apply Lemma~\ref{lem:quadraticelimination}. This estimate is
independent of the degrees of $B,g$ and the dimension $s$.
\end{proof}

\subsubsection*{A nonisotropic repeated factor}

Normalize the repeated factor to $u$ with $\E[u^2]=1$ and take
an orthogonal Gaussian complement $Y\in\C^2$. Every such cubic
has the form in Lemma~\ref{lem:quadraticfibrehighdegree}, with
\[
 B(u)=p u^2+q u+r,\qquad\deg g\le3.
\]
For a double line $p\ne0$. If its effective quartic coefficient
is nonzero, the lemma excludes purity immediately. If that
coefficient vanishes identically, then
\[
 0=p^T\operatorname{adj}(I-2tA)p
   =p^Tp-2t\bigl(\operatorname{tr}(A)p^Tp-p^TAp\bigr).
\]
Thus $p^Tp=p^TAp=0$. The nonzero vector $(0,p)$ is isotropic,
lies in the leading singular plane $u=0$, and satisfies
$F_2(0,p)=0$. It supplies the frame of Lemma~\ref{lem:frame}.
For a triple line, $p=0$ and $[u^3]g\ne0$. The effective
phase has exact cubic coefficient $t[u^3]g$, so the same
high-degree lemma excludes purity. This exhausts all matrices
$A$, including nilpotent and defective ones, without a matrix
case split.

\subsubsection*{An isotropic repeated factor}

Use Gaussian circular coordinates with repeated factor $Z$.
The singular plane is $Z=0$, whose only isotropic direction is
the $W$ axis. Write $\lambda=[W^2]F_2$. If $\lambda=0$, the
frame test applies. Assume $\lambda\ne0$.

For a double line write its other factor as
$aW+bT+cZ$. The Gaussian null rotation
\[
 Z=Z',\quad T=T'+sZ',\quad W=W'-sT'-s^2Z'/2
\]
removes the $T$ term if $a\ne0$, or removes the $Z$ term if
$a=0,b\ne0$. Scaling $P$ gives exactly these leading forms:
\[
 \mathrm I: Z^2W+a_0Z^3,\qquad
 \mathrm {II}:Z^2T,\qquad \mathrm {III}:Z^3.
\]
The last is the triple line. The rotation preserves the Gaussian
and keeps the $W^2$ coefficient nonzero. All lower terms remain
arbitrary.

Write $P=\lambda W^2+W H(Z,T)+G(Z,T)$. In type I,
$H=Z^2+bZ+cT+d$ and $G=a_0Z^3+G_2$. In types II and III,
$H=bZ+cT+d$ and $G=Z^2T+G_2$ or $Z^3+G_2$ respectively.
Here $G_2$ is an arbitrary polynomial of degree at most two.
Eliminating $W$ in $tP-ZW-T^2/2$ gives
\[
 \Psi=t\left(G-\frac{H^2}{4\lambda}\right)
       +\frac{ZH}{2\lambda}-\frac{T^2}{2}
       -\frac{Z^2}{4\lambda t}.
\]
Its $T^2$ coefficient is $C_T=a t-1/2$, where
$a=g_{02}-c^2/(4\lambda)$ and $g_{02}=[T^2]G_2$.
Eliminate $T$ as well. If $e=\deg_t C_T\in\{0,1\}$, the total
normalization has $v_\infty(h)=(1+e)/2$. The actual remaining
leading coefficients are
\[
 \begin{array}{c|c|c}
 \text{type}&[Z^d]\varphi&d\\\hline
 \mathrm I&\displaystyle
  -\frac{t}{4\lambda}-\frac{c^2t^2}{16\lambda^2C_T}
   =-\frac{t(2g_{02}t-1)}{8\lambda C_T}&4\\[5pt]
 \mathrm {II}&\displaystyle-\frac{t^2}{4C_T}&4\\[4pt]
 \mathrm {III}&t&3.
 \end{array}
\]
None is identically zero. Their valuations are at most $e-1$,
$e-2$, and $-1$, respectively. Each is strictly less than
$d(1+e)/2$. Lemma~\ref{lem:quadraticelimination} excludes
purity in all three types. In particular this includes the
boundary where the quartic coefficient tends to a nonzero
constant while lower coefficients grow.

\begin{proposition}\label{prop:repeatedcomplete}
Every pure ternary cubic with repeated leading factor has a
Gaussian affine-linear coordinate. It therefore satisfies
$m>3\deg Q$.
\end{proposition}
\begin{proof}
The nonisotropic and isotropic alternatives above exhaust the
repeated factor. In every surviving branch the explicit frame
test succeeds. Apply the complete affine-linear theorem.
\end{proof}


\subsection{The geometric case split for frame extraction}

For a cubic of degree exactly three, either the frame test
\eqref{eq:frame} succeeds, or it does not. In the latter case,
the projective curve $F_3=0$ belongs to exactly one of these types:
\begin{center}\small
\begin{tabular}{@{}p{0.33\textwidth}p{0.56\textwidth}@{}}
\toprule
Leading curve & No-frame exclusion\\
\midrule
Smooth cubic & Corollary~\ref{cor:smoothleading}\\
Irreducible nodal cubic & Corollary~\ref{cor:ordinarycomplete}\\
Irreducible cuspidal cubic & Proposition~\ref{prop:cusptangentcomplete}\\
Line and smooth conic, transverse & Corollary~\ref{cor:ordinarycomplete}\\
Line and smooth conic, tangent & Proposition~\ref{prop:cusptangentcomplete}\\
Three nonconcurrent lines & Corollary~\ref{cor:ordinarycomplete}\\
Three concurrent distinct lines & Proposition~\ref{prop:concurrentcomplete}\\
Double line and a distinct line & Proposition~\ref{prop:repeatedcomplete}\\
Triple line & Proposition~\ref{prop:repeatedcomplete}\\
\bottomrule
\end{tabular}
\end{center}
The list is exhaustive by factorization of a plane cubic over $\C$.
Every row has now been proved, including its lower terms and all
critical degeneracies. In the repeated-line rows the singular set
is a whole projective line; the proof checks its isotropic directions
rather than choosing an arbitrary point on it.

\begin{proof}[Proof of Proposition~\ref{prop:frameextract}]
Suppose the frame test fails for a pure polynomial of degree three.
Its leading curve is one of the nine types in the table. The smooth
case is excluded by the gradient criterion, and the ordinary-node
cases by Corollary~\ref{cor:ordinarycomplete}. Propositions
\ref{prop:cusptangentcomplete}, \ref{prop:concurrentcomplete} and
\ref{prop:repeatedcomplete} force a frame in every other case.
All possibilities give a contradiction. This proves frame extraction.
\end{proof}
